% Modified 31 Oct 2005:  Conditioning fallacy alluded to.
% This chapter has been modified on 6-4-05.
% There are two \choice
\pagestyle{headings}
\setcounter{chapter}{0}
\chapter{Distribusi Peluang Diskret} \label{chp 1}
\setcounter{page}{1}

\section{Simulasi Peluang Diskret} \label{sec 1.1}

\subsection*{Peluang}

Dalam bab ini, pertama-tama kita akan membahas percobaan acak dengan sejumlah
terhingga hasil yang mungkin, yaitu $\omega_1$, $\omega_2$, \dots, $\omega_n$.  Sebagai contoh, kita
melempar sebuah dadu dan hasil yang mungkin adalah 1, 2, 3, 4, 5, 6, sesuai dengan
sisi yang muncul.  Kita melempar sebuah koin dengan hasil yang mungkin H (kepala) dan T
(ekor).  
\par
Sering kali kita perlu dapat merujuk pada suatu hasil percobaan.  Sebagai 
contoh, kita mungkin ingin menuliskan ekspresi matematika yang menyatakan jumlah dari empat kali pelemparan
sebuah dadu.  Untuk itu, kita dapat menetapkan $X_i$, $i = 1, 2, 3, 4,$ sebagai nilai
hasil dari keempat pelemparan tersebut, lalu menuliskan ekspresi
$$X_1 + X_2 + X_3 + X_4$$
untuk jumlah keempat pelemparan.  Besaran $X_i$ disebut \emx {peubah acak}\index{peubah 
acak}.
Peubah acak hanyalah sebuah ekspresi yang nilainya merupakan hasil suatu
percobaan tertentu.  Seperti jenis peubah lain dalam matematika, peubah acak
dapat memiliki nilai yang berbeda-beda.
\par
Misalkan $X$ adalah peubah acak yang menyatakan hasil pelemparan sebuah dadu.  Kita akan menetapkan
peluang bagi setiap hasil yang mungkin dari percobaan ini.  Caranya adalah dengan memasangkan pada setiap
hasil $\omega_j$ sebuah bilangan taknegatif $m(\omega_j)$ sedemikian rupa sehingga
$$ 
m(\omega_1) + m(\omega_2) + \cdots + m(\omega_6) = 1\ .
$$
Fungsi $m(\omega_j)$ disebut \emx {fungsi distribusi}\index{fungsi
distribusi} dari peubah acak
$X$.  Dalam kasus pelemparan dadu, kita menetapkan peluang yang sama, yakni
peluang 1/6, pada setiap hasil.  Dengan penetapan peluang ini, kita dapat menulis
$$
P(X \le 4) = {2\over 3}
$$
untuk menyatakan bahwa peluangnya adalah $2/3$ bagi sebuah pelemparan dadu untuk menghasilkan nilai yang
tidak melebihi 4.  
\par
Misalkan $Y$ adalah peubah acak yang menyatakan hasil pelemparan sebuah koin.  Dalam kasus ini, terdapat
dua hasil yang mungkin, yang dapat kita beri label H dan T.  Kecuali jika ada alasan untuk menduga bahwa
salah satu sisi koin muncul lebih sering daripada sisi lainnya, wajar jika kita menetapkan
peluang 1/2 pada masing-masing hasil.
\par
Dalam kedua percobaan di atas, setiap hasil diberi peluang yang sama.  Tentu saja, secara umum
tidak selalu demikian.  Sebagai contoh, jika suatu obat diketahui efektif pada 30 persen dari seluruh
penggunaannya, kita dapat menetapkan peluang .3 bahwa obat itu efektif pada penggunaan berikutnya dan
.7 bahwa obat itu tidak efektif.  Contoh terakhir ini menggambarkan \emx {konsep frekuensi
peluang}\index{peluang!konsep frekuensi}.  Artinya, jika kita memiliki peluang $p$ bahwa 
suatu percobaan memberikan hasil $A$, maka apabila percobaan itu diulangi berkali-kali, 
kita mengharapkan bahwa
proporsi kemunculan $A$ kira-kira sebesar $p$.  Untuk menguji gagasan intuitif semacam ini, kita akan
terbantu dengan menelaah beberapa masalah tersebut melalui percobaan.  Sebagai 
contoh, kita dapat melempar koin berkali-kali dan melihat apakah proporsi kemunculan sisi kepala 
mendekati 1/2.  Kita juga dapat menyimulasikan percobaan ini pada komputer.

\subsection*{Simulasi}

Kita ingin dapat melakukan percobaan yang bersesuaian dengan suatu himpunan
peluang tertentu; misalnya, $m(\omega_1) = 1/2$, $m(\omega_2) = 1/3$, dan
$m(\omega_3) = 1/6$.  Dalam hal ini, tiga sisi sebuah dadu bersisi enam dapat diberi tanda $\omega_1$,
dua sisi diberi tanda $\omega_2$, dan satu sisi diberi tanda $\omega_3$.  
\par
Dalam kasus umum, kita mengasumsikan bahwa $m(\omega_1)$, $m(\omega_2)$, \dots, $m(\omega_n)$ semuanya merupakan bilangan
rasional dengan penyebut bersama terkecil $n$.  Jika $n > 2$, kita dapat membayangkan sebuah dadu silindris panjang dengan
penampang berupa segi-$n$ beraturan.  Jika $m(\omega_j) = n_j/n$, kita dapat memberi tanda pada $n_j$ dari
sisi panjang silinder itu dengan $\omega_j$; jika salah satu sisi ujung yang muncul, kita cukup
melempar dadu tersebut lagi.  Jika $n = 2$, percobaan dapat dilakukan dengan sebuah koin.
\par
Kita terutama akan menaruh perhatian pada pengulangan suatu percobaan acak
berkali-kali.  Meskipun dadu silindris memudahkan kita melakukan
beberapa pengulangan, cara itu akan sulit digunakan untuk menjalankan sangat banyak
percobaan.  Karena komputer modern dapat melakukan banyak operasi
dalam waktu yang sangat singkat, wajar jika kita menggunakan komputer untuk tugas ini.

\subsection*{Bilangan Acak}

Pertama-tama kita harus mencari padanan pelemparan dadu pada komputer.  Hal ini dilakukan
dengan menggunakan \emx {generator bilangan acak}.\index{generator bilangan acak}
Bergantung pada paket perangkat lunak yang digunakan,
komputer dapat diminta menghasilkan bilangan real antara 0 dan 1, atau bilangan bulat dalam suatu
himpunan bilangan bulat berurutan.  Dalam kasus pertama, bilangan real dipilih sedemikian rupa sehingga
peluang bilangan tersebut terletak pada sembarang subinterval tertentu dari interval satuan ini
sama dengan panjang subinterval tersebut.  Dalam kasus kedua, setiap bilangan bulat memiliki
peluang yang sama untuk dipilih.
\par
Misalkan $X$ adalah peubah acak dengan fungsi distribusi $m(\omega)$, dengan $\omega$ berada dalam himpunan $\{\omega_1,
\omega_2, \omega_3\}$, serta $m(\omega_1) = 1/2$, $m(\omega_2) = 1/3$, dan $m(\omega_3) = 1/6$.  Jika
paket komputer kita dapat menghasilkan bilangan bulat acak dalam himpunan $\{1, 2, ..., 6\}$, kita cukup memintanya
melakukan hal itu, lalu mengaitkan 1, 2, dan 3 dengan $\omega_1$, 4 dan 5 dengan $\omega_2$, serta 6
dengan $\omega_3$.  Jika paket komputer kita menghasilkan bilangan real acak $r$ dalam interval
$(0,~1)$, maka ekspresi
$$\lfloor {6r}\rfloor + 1$$
akan menjadi bilangan bulat acak antara 1 dan 6.  (Notasi $\lfloor x \rfloor$ berarti bilangan bulat terbesar
yang tidak melebihi $x$, dan dibaca ``lantai dari $x$.")
\par
Metode yang digunakan komputer untuk menghasilkan bilangan real acak dijelaskan dalam pembahasan
sejarah pada akhir bagian ini.  Contoh berikut memberikan cuplikan keluaran program {\bf
RandomNumbers}.\index{RandomNumbers (program)}
\begin{example}(Pembangkitan Bilangan Acak) \label{exam 1.05}
Program {\bf RandomNumbers} menghasilkan $n$ bilangan real acak dalam interval $[0, 1]$, dengan
$n$ dipilih oleh pengguna.  Ketika kami menjalankan program dengan $n = 20$, kami memperoleh data yang ditampilkan pada
Tabel~\ref{table 1.1}.
\begin{table}
\centering
\begin{tabular}{llll} 
.203309 &
.762057 &
.151121 &
.623868 \\
.932052 &
.415178 &
.716719 &
.967412 \\
.069664 & 
.670982 &
.352320 &
.049723 \\
.750216 &
.784810 &
.089734 &
.966730 \\
.946708 &
.380365 &
.027381 &
.900794 \\
\end{tabular}
\caption{Contoh keluaran program {\bf RandomNumbers}.}
\label{table 1.1}
\end{table}
\end{example}

\begin{example}(Pelemparan Koin) \label{exam 1.1}
 Seperti telah kita catat, intuisi kita menyatakan bahwa peluang
memperoleh sisi kepala dalam sekali pelemparan koin adalah 1/2.  Agar komputer
melempar koin, kita dapat memintanya memilih bilangan real acak dalam interval $[0, 1]$, lalu
memeriksa apakah bilangan itu kurang dari 1/2.  Jika ya, hasilnya kita sebut \emx {kepala}; jika
tidak,
kita menyebutnya \emx {ekor}.  Cara lain adalah meminta komputer memilih bilangan bulat
acak dari himpunan $\{0, 1\}$.  Program {\bf CoinTosses}\index{CoinTosses (program)} menjalankan
percobaan pelemparan koin sebanyak $n$ kali.  Menjalankan program ini dengan
$n = 20$ menghasilkan:
\vskip .1in
\begin{center}
THTTTHTTTTHTTTTTHHTT.
\end{center}
\vskip .1in

Perhatikan bahwa dalam 20 pelemparan, kita memperoleh 5 kepala dan 15 ekor.  Mari kita lempar
koin sebanyak $n$ kali, dengan $n$ jauh lebih besar dari 20, lalu kita lihat apakah proporsi kepala yang diperoleh lebih dekat dengan
dugaan intuitif kita, yaitu 1/2.  Program {\bf CoinTosses} mencatat jumlah kepala.  Ketika
kita menjalankan program ini dengan $n = 1000$, kita memperoleh 494 kepala.  Ketika kita menjalankannya dengan $n = 10000$, kita memperoleh
5039 kepala.

Kita melihat bahwa ketika koin dilempar 10{,}000 kali, proporsi
kepala
mendekati ``nilai sebenarnya" .5 untuk memperoleh sisi kepala ketika sebuah koin
dilempar.  Model matematika untuk percobaan ini disebut Percobaan Bernoulli (lihat Bab~
\ref{chp 3}).
\emx {Hukum Bilangan Besar}, yang akan kita pelajari kemudian (lihat Bab~\ref{chp
8}), 
akan menunjukkan bahwa dalam model Percobaan Bernoulli, proporsi kepala seharusnya mendekati .5,
sesuai dengan gagasan intuitif kita mengenai interpretasi frekuensi dari peluang.
\par
Tentu saja, program kita dapat dengan mudah diubah untuk menyimulasikan koin
yang peluang munculnya sisi kepala adalah $p$, dengan $p$ merupakan bilangan real antara 0 dan 1.
\end{example}

\par
Dalam kasus pelemparan koin, kita telah mengetahui peluang terjadinya kejadian
pada setiap percobaan.  Kekuatan simulasi yang sesungguhnya terletak pada
kemampuannya memperkirakan peluang yang belum diketahui sebelumnya. 
Metode
ini telah digunakan dalam penemuan strategi-strategi baru yang membuat permainan kasino
blackjack menguntungkan pemain.  Kita akan menggambarkan gagasan ini melalui
sebuah
situasi sederhana yang memungkinkan kita menghitung peluang sebenarnya dan melihat seberapa
efektif simulasi tersebut.

\begin{example}(Pelemparan Dadu)\label{exam 1.2}
Kita membahas sebuah permainan dadu yang berperan penting dalam
perkembangan sejarah peluang.  Surat-surat terkenal
antara Pascal\index{PASCAL, B.} dan Fermat\index{FERMAT, P.}, 
yang oleh banyak orang dianggap mengawali kajian serius tentang
peluang, dipicu oleh permintaan bantuan dari seorang bangsawan sekaligus
penjudi Prancis, Chevalier de~M\'er\'e\index{de M\'ER\'E, CHEVALIER}.  
Konon, de~M\'er\'e bertaruh
bahwa dalam empat kali pelemparan sebuah dadu, setidaknya satu angka enam akan muncul.  Ia menang
secara konsisten dan, untuk menarik lebih banyak orang agar bermain, ia mengubah permainan menjadi taruhan
bahwa
dalam 24 kali pelemparan dua dadu, sepasang angka enam akan muncul. Dikisahkan bahwa
de~M\'er\'e kalah dengan 24 pelemparan dan merasa bahwa diperlukan 25 pelemparan
agar
permainan itu menguntungkan.  Bahwa matematika ternyata keliru merupakan \emx {un grand scandale}.
\par
Kita akan mencoba melihat apakah de~M\'er\'e benar dengan menyimulasikan berbagai taruhannya. 
Program {\bf DeMere1}\index{DeMere1 (program)} menyimulasikan sejumlah besar percobaan 
dan, dalam setiap percobaan, memeriksa apakah angka enam muncul dalam empat kali pelemparan dadu. Ketika program ini dijalankan untuk
1000 permainan, angka enam muncul dalam empat pelemparan pertama pada 48.6 persen permainan.  Ketika program dijalankan
untuk 10,000 permainan, hal ini terjadi pada 51.98 persen permainan.
\par
Kita mencatat bahwa hasil pelaksanaan kedua menunjukkan bahwa de~M\'er\'e benar
ketika meyakini bahwa taruhannya dengan satu dadu menguntungkan; namun, seandainya
kesimpulan kita didasarkan pada pelaksanaan pertama, kita akan memutuskan bahwa ia keliru. 
\emx {Hasil simulasi yang akurat memerlukan sejumlah besar percobaan.}
\end{example}
\par
Program {\bf DeMere2}\index{DeMere2 (program)} menyimulasikan taruhan kedua de~M\'er\'e, yaitu bahwa 
sepasang angka enam akan muncul dalam 
$n$ kali pelemparan sepasang dadu.  Simulasi sebelumnya menunjukkan bahwa penting untuk mengetahui berapa
banyak uji coba yang harus kita simulasikan agar dapat mengharapkan tingkat ketelitian tertentu dalam
pendekatan kita.  Kelak akan kita lihat bahwa dalam percobaan jenis
ini, aturan praktis kasarnya adalah bahwa, setidaknya pada 95\% kesempatan, galat tidak melebihi
kebalikan dari akar kuadrat jumlah uji coba.  Untungnya, untuk permainan dadu ini, 
peluang eksaknya mudah dihitung.  Pada bagian berikutnya akan kita tunjukkan bahwa untuk taruhan pertama,
peluang de~M\'er\'e menang adalah $1 - (5/6)^4 = .518$. 
\par 
Perhitungan ini dapat dipahami sebagai berikut:  Peluang angka 6 tidak
muncul pada pelemparan pertama adalah $(5/6)$.  Peluang angka 6 tidak muncul pada
kedua pelemparan pertama adalah $(5/6)^2$.  Dengan penalaran yang sama,
peluang angka 6 tidak muncul pada satu pun dari empat pelemparan pertama adalah $(5/6)^4$. 
Jadi, peluang setidaknya satu angka 6 muncul dalam empat pelemparan pertama adalah $1 -
(5/6)^4$.  Demikian pula, untuk taruhan kedua dengan 24 pelemparan, peluang
de~M\'er\'e
menang adalah $1 - (35/36)^{24} = .491$, sedangkan untuk 25 pelemparan peluangnya adalah $1 - (35/36)^{25}
= .506$.  
\par
Menurut aturan praktis yang disebutkan di atas,
diperlukan 27,000 pelemparan agar terdapat kemungkinan yang wajar untuk menentukan peluang-peluang ini dengan
ketelitian yang cukup guna memastikan bahwa keduanya berada pada sisi yang berlawanan dari .5.  Menarik untuk direnungkan
apakah seorang penjudi dapat mengenali peluang seperti itu dengan ketelitian yang diperlukan hanya dari pengalaman
berjudi.  Beberapa penulis sejarah peluang menduga bahwa de~M\'er\'e sebenarnya
hanya tertarik pada persoalan-persoalan tersebut sebagai masalah peluang yang memikat.

\begin{example}(Kepala atau Ekor)\label{exam 1.3}
Untuk contoh berikutnya, kita membahas suatu masalah
yang jawaban eksaknya sulit diperoleh, tetapi hasil kualitatifnya mudah
diperoleh melalui simulasi.  Peter dan Paul memainkan permainan yang disebut \emx {kepala
atau
ekor}.  Dalam permainan ini, sebuah koin seimbang dilempar berulang kali---kita pilih
40 kali.  Setiap kali sisi kepala muncul, Peter memenangkan 1 penny dari Paul, dan setiap kali sisi
ekor muncul, Peter kehilangan 1 penny kepada Paul.  Sebagai contoh, jika hasil
40 pelemparan tersebut adalah
\begin{center}
THTHHHHTTHTHHTTHHTTTTHHHTHHTHHHTHHHTTTHH.
\end{center}
\noindent Perolehan Peter dapat digambarkan seperti pada Gambar~\ref{fig 1.3}.
\putfig{3.5truein}{PSfig1-3}{Perolehan Peter dalam 40 permainan kepala atau ekor.}{fig 1.3}
\par
Dalam permainan ini, Peter memenangkan 6 penny.  Wajar jika kita menanyakan
peluang bahwa ia akan memenangkan $j$ penny; di sini $j$ dapat berupa bilangan genap
mana pun dari $-40$ sampai $40$.  Masuk akal untuk menduga bahwa nilai $j$ dengan
peluang tertinggi adalah $j = 0$, sebab hal ini terjadi ketika jumlah kepala
sama dengan jumlah ekor.  Demikian pula, kita dapat menduga bahwa nilai $j$ dengan
peluang terendah adalah $j = \pm 40$.
\par
Pertanyaan menarik kedua tentang permainan ini adalah: Berapa kali
dalam 40 pelemparan Peter akan berada dalam posisi unggul?  Dengan melihat grafik
perolehannya (Gambar~\ref{fig 1.3}), kita melihat bahwa Peter unggul ketika perolehannya 
positif. Namun,
kita harus menetapkan suatu kesepakatan ketika perolehannya 0 jika kita ingin semua pelemparan
turut dihitung dalam jumlah saat ia unggul.  Kita menggunakan kesepakatan
bahwa ketika perolehan Peter 0, ia dianggap unggul jika ia unggul pada
pelemparan sebelumnya, dan tidak dianggap unggul jika sebelumnya ia tertinggal.  Dengan
kesepakatan ini, Peter unggul 34 kali dalam contoh kita.  Sekali lagi,
intuisi kita mungkin menyatakan bahwa jumlah saat unggul yang paling mungkin adalah
1/2 dari 40, yaitu 20, sedangkan jumlah yang paling kecil kemungkinannya adalah kasus ekstrem 40
atau 0.
\par
Hal ini mudah dipastikan dengan menyimulasikan permainan berkali-kali dan
mencatat berapa kali perolehan akhir Peter bernilai $j$, serta
berapa kali Peter akhirnya unggul sebanyak $k$ kali.  Proporsi
dari seluruh permainan kemudian memberikan perkiraan bagi
peluang yang bersesuaian.  Program {\bf HTSimulation}\index{HTSimulation (program)} menjalankan
simulasi ini.  Perhatikan bahwa jika jumlah pelemparan dalam permainan genap,
jumlah saat berada dalam posisi unggul hanya mungkin berupa bilangan genap.  Kita telah menyimulasikan 
permainan ini 10{,}000 kali. 
Hasilnya ditampilkan pada Gambar~\ref{fig 1.4.5} dan \ref{fig 1.4.6}.  Grafik-grafik ini,
yang kita sebut grafik paku\index{grafik paku}, dibuat menggunakan program {\bf Spikegraph}.
\index{Spikegraph (program)}  
Garis vertikal, atau paku, pada posisi $x$ di sumbu horizontal memiliki tinggi yang sama dengan
proporsi hasil yang bernilai
$x$.  Intuisi kita mengenai perolehan akhir Peter ternyata cukup tepat, tetapi intuisi kita
mengenai jumlah saat Peter unggul sama sekali keliru.  Simulasi
menunjukkan bahwa jumlah saat unggul yang paling kecil kemungkinannya adalah 20, sedangkan yang paling mungkin adalah 0
atau 40.  Hal ini memang benar, dan penjelasannya dapat dibayangkan dengan memainkan
permainan kepala atau ekor dengan sangat banyak pelemparan dan mengamati grafik
perolehan Peter.  Pada Gambar~\ref{fig 1.4.1} kita menampilkan hasil simulasi permainan untuk
1000 pelemparan, dan pada Gambar~\ref{fig 1.4.2} untuk 10{,}000 pelemparan.
\par
Dalam contoh kedua, Peter unggul hampir sepanjang waktu.  Meskipun demikian, terdapat fakta yang luar biasa:
jika permainan dilanjutkan cukup lama, perolehan Peter akan terus
kembali ke 0, tetapi terdapat selang waktu yang sangat panjang antara dua kejadian
semacam itu.  Hasil ini dan hasil-hasil terkait akan dibahas dalam
Bab~\ref{chp 12}.
\end{example}

Dalam semua contoh sejauh ini, kita telah menyimulasikan hasil-hasil yang berpeluang sama.  Selanjutnya,
kita akan melihat contoh dengan hasil-hasil yang peluangnya tidak sama.

\putfig{4truein }{PSfig1-4-5}{Distribusi perolehan.}{fig 1.4.5}

\putfig{4truein }{PSfig1-4-6}{Distribusi jumlah saat berada dalam posisi unggul.}{fig 1.4.6}

\putfig{4truein}{PSfig1-4}{Perolehan Peter dalam 1000 permainan kepala atau ekor.}{fig 1.4.1}

\putfig{4truein} {PSfig1-5}{Perolehan Peter dalam 10,000 permainan kepala atau ekor.}{fig 1.4.2}


\begin{example}(Pacuan Kuda)\label{exam 1.4}
Empat ekor kuda (Acorn, Balky, Chestnut, dan Dolby) telah berpacu berkali-kali.
Diperkirakan bahwa Acorn menang pada 30 persen perlombaan, Balky 40
persen, Chestnut 20 persen, dan Dolby 10 persen.
\par
Kita dapat meminta komputer menyelenggarakan satu perlombaan sebagai berikut:  
Pilih bilangan acak $x$.  Jika $x < .3$, kita katakan Acorn menang.  Jika $.3 \le x < .7$,
Balky menang.  Jika $.7 \le x < .9$, Chestnut menang.  Terakhir, jika $.9 \le x$,
Dolby menang. 
\par
Program {\bf HorseRace}\index{HorseRace (program)} menggunakan metode ini untuk menyimulasikan hasil
$n$ perlombaan.   Setelah menjalankan program ini untuk $n = 10$, kita mendapati bahwa Acorn menang pada 40 persen perlombaan,
Balky 20 persen, Chestnut 10 persen, dan Dolby 30 persen.  
Diperlukan lebih banyak perlombaan agar hasilnya lebih sesuai dengan pengalaman sebelumnya.
Karena itu, kita menjalankan program untuk menyimulasikan 1000 perlombaan dengan keempat kuda kita. 
Meskipun sangat lelah setelah semua perlombaan ini, penampilan mereka cukup sesuai
dengan perkiraan kemampuan masing-masing. Acorn menang pada 29.8 persen perlombaan,
Balky 39.4 persen, Chestnut 19.5 persen, dan Dolby 11.3 persen. 
\par
Program {\bf GeneralSimulation}\index{GeneralSimulation (program)} menggunakan metode ini untuk menyimulasikan
pengulangan sembarang percobaan dengan sejumlah terhingga hasil yang memiliki
peluang diketahui.
\end{example}

\subsection*{Catatan Sejarah}

Siapa pun yang memainkan permainan acak yang sama berulang-ulang sesungguhnya sedang menjalankan
simulasi; dalam pengertian ini, proses simulasi telah berlangsung selama
berabad-abad.  Seperti telah kita catat, banyak masalah awal dalam kajian peluang
mungkin saja terilhami oleh pengalaman para penjudi.  
\par
Wajar bagi siapa pun yang mencoba memahami teori peluang untuk melakukan
percobaan sederhana dengan melempar koin, melempar dadu, dan sebagainya.  Naturalis
Buffon\index{BUFFON, G. L.} melempar sebuah koin 4040 kali dan memperoleh 2048 kepala serta 1992 ekor.  Ia
juga memperkirakan bilangan $\pi$ dengan menjatuhkan jarum pada permukaan bergaris dan
mencatat berapa kali jarum melintasi sebuah garis\choice{.\footnote{Lihat Bagian 2.1
dalam buku lengkap Grinstead-Snell. Hal ini juga diperagakan dalam Probability
Demo.}}{\ (lihat Bagian~\ref{sec 2.1}).}   Biolog Inggris W. F. R. Weldon\index{WELDON,
W. F. R.}\footnote{T.~C.~Fry,
\emx {Probability and Its Engineering Uses,} 2nd~ed. (Princeton: Van Nostrand, 1965).}
mencatat 26{,}306 kali pelemparan 12 dadu, sedangkan
ilmuwan Swiss Rudolf Wolf\index{WOLF, R.}\footnote{E.~Czuber, \emx
{Wahrscheinlichkeitsrechnung,} 3rd~ed. (Berlin: Teubner, 1914).} mencatat
100{,}000 kali pelemparan sebuah dadu tanpa komputer.  Percobaan semacam itu
sangat
memakan waktu dan mungkin tidak menggambarkan secara akurat fenomena acak yang sedang
dikaji.  Sebagai contoh, analisis lebih lanjut atas data yang dicatat dalam
percobaan dadu Weldon dan Wolf menimbulkan dugaan bahwa dadu mereka tidak seimbang.
Statistikawan Karl Pearson\index{PEARSON, K.} menganalisis sejumlah besar hasil pada beberapa
meja rolet dan menduga bahwa roda-roda tersebut tidak seimbang.  Pada 1894 ia menulis:

\begin{quote} 
Jelaslah, karena Kasino tidak menjalankan tujuan berharga sebagai sebuah
laboratorium besar untuk penyusunan statistik peluang, kasino tidak memiliki
\emx {raison d'\^etre} ilmiah.  Para ilmuwan tidak boleh membiarkan teori-teori mereka yang paling
halus diabaikan dengan cara yang tidak tahu malu ini!  Pemerintah Prancis harus
didesak oleh hierarki ilmu pengetahuan untuk menutup ruang-ruang perjudian; tentu
saja, akan merupakan tindakan yang mulia jika sumber daya Kasino yang tersisa diserahkan kepada
Acad\'emie des Sciences untuk mendanai sebuah laboratorium peluang
ortodoks; terutama bagi cabang baru kajian tersebut, yakni penerapan
teori peluang pada masalah-masalah biologis evolusi, yang tampaknya akan
menyita begitu banyak pemikiran manusia dalam waktu dekat.\footnote{K. Pearson,
``Science and Monte Carlo," \emx {Fortnightly Review}, vol.~55 (1894),
p.~193;
disitir dalam S. M. Stigler, \emx {The History of Statistics} (Cambridge: Harvard
University Press, 1986).}
\end{quote}

Meskipun demikian, percobaan-percobaan awal ini memberikan petunjuk dan menghasilkan
penemuan penting dalam peluang dan statistika.  Percobaan tersebut mengantarkan Pearson pada \emx
{uji khi-kuadrat}, yang sangat penting untuk menguji apakah data
yang diamati sesuai dengan suatu distribusi peluang tertentu.

Pada awal 1900-an, jelas bahwa diperlukan cara yang lebih baik untuk menghasilkan bilangan acak.
Pada 1927, L. H. C. Tippett\index{TIPPETT, L. H. C.} menerbitkan daftar 41{,}600 digit
yang diperoleh dengan memilih bilangan secara sembarang dari laporan sensus.  Pada 1955, RAND
Corporation\index{RAND Corporation} mencetak tabel berisi 1{,}000{,}000 bilangan acak yang dihasilkan dari
derau elektronik.  Kehadiran komputer berkecepatan tinggi membuka
kemungkinan
untuk menghasilkan bilangan acak secara langsung pada komputer, dan pada akhir 1940-an
John von~Neumann\index{von NEUMANN, J.} mengusulkan cara berikut: Misalkan Anda menginginkan
barisan acak bilangan empat digit.  Pilih sembarang bilangan empat digit, misalnya
6235, sebagai awal.  Kuadratkan bilangan ini untuk memperoleh 38{,}875{,}225.  Sebagai bilangan
kedua, pilih empat digit tengah dari kuadrat tersebut (yaitu 8752).  Ulangi
proses yang sama mulai dari 8752 untuk memperoleh bilangan ketiga, dan seterusnya.  
\par
Metode yang lebih modern melibatkan
konsep aritmetika modular\index{aritmetika modular}.  
Jika $a$ merupakan bilangan bulat dan $m$ merupakan bilangan bulat positif, maka
$a\ (\mbox{mod}\  m)$
berarti sisa pembagian $a$ oleh $m$.  Sebagai contoh, $10\ (
\mbox{mod}\  4) = 2$,
$8\ (\mbox{mod}\  2) = 0$, dan seterusnya.  Untuk menghasilkan barisan acak $X_0,
X_1, X_2,
\dots$, pilih bilangan awal $X_0$, lalu peroleh bilangan
$X_{n+1}$ dari $X_n$ dengan rumus 
$$X_{n+1} = (aX_n + c)\ (\mbox{mod}\ m)\ ,$$
di mana $a$, $c$, dan $m$ merupakan konstanta yang dipilih dengan cermat.  Barisan 
$ X_0,  X_1,$ 
$X_2, \dots$ 
kemudian menjadi barisan bilangan bulat antara 0 dan $m-1$.  Untuk memperoleh barisan bilangan real dalam $[0,1)$,
kita membagi setiap $X_j$ dengan $m$.  Barisan yang dihasilkan terdiri atas bilangan rasional berbentuk $j/m$,
dengan $0 \le j \le m-1$.  Karena $m$ biasanya merupakan bilangan bulat yang sangat besar, kita menganggap bilangan dalam barisan
tersebut sebagai bilangan real acak dalam $[0, 1)$.
\par
Baik dalam metode penguadratan von~Neumann maupun teknik aritmetika modular, barisan
bilangan sebenarnya sepenuhnya ditentukan oleh bilangan pertama.  Jadi,
tidak ada sesuatu yang sungguh-sungguh acak dalam barisan ini.  Meskipun demikian, barisan tersebut menghasilkan
bilangan yang perilakunya sangat menyerupai prediksi teori bagi percobaan acak.  Untuk
memperoleh barisan yang berbeda bagi percobaan yang berbeda, bilangan awal $X_0$
dipilih dengan prosedur lain yang, misalnya, dapat melibatkan waktu
saat itu.\footnote{Untuk pembahasan terperinci tentang bilangan acak, lihat D.~E.~Knuth, \emx {The Art of
Computer Programming,} vol.~II (Reading: Addison-Wesley, 1969).}
\par
Selama Perang Dunia Kedua, para fisikawan di Los Alamos Scientific
Laboratory
perlu mengetahui, untuk keperluan perisai radiasi, seberapa jauh neutron bergerak menembus
berbagai bahan.  Pertanyaan ini berada di luar jangkauan perhitungan
teoretis. Daniel McCracken\index{McCRACKEN, D.}, dalam tulisannya di \emx
{Scientific American}, merangkum data fisik yang tersedia bagi para peneliti:
jarak rata-rata yang ditempuh neutron berkecepatan tertentu sebelum bertumbukan
dengan inti atom, peluang neutron dipantulkan alih-alih diserap, dan sebaran
energi yang hilang dalam tumbukan.\footnote{D.~D.~McCracken, ``The Monte Carlo
Method," \emx {Scientific American,} vol.~192 (May 1955), p.~90. Ringkasan
editorial; ungkapan sumber tidak diterjemahkan atau direproduksi.}

John von~Neumann\index{von NEUMANN, J.} dan Stanislas Ulam\index{ULAM, S.} 
mengusulkan agar masalah tersebut diselesaikan dengan
memodelkan percobaan menggunakan perangkat acak pada komputer.  Karena pekerjaan mereka
bersifat rahasia, pekerjaan itu perlu diberi nama sandi.  Von~Neumann memilih nama
``Monte Carlo."  Sejak saat itu, metode simulasi ini disebut
\emx {Metode Monte Carlo}.

William Feller\index{FELLER, W.} menunjukkan kemungkinan penggunaan simulasi
komputer untuk menggambarkan konsep-konsep dasar peluang dalam bukunya \emx
{An Introduction to Probability Theory and Its Applications.}  Ketika membahas
masalah jumlah saat unggul dalam permainan ``kepala atau ekor," Feller
menekankan bahwa fluktuasi yang diprediksi teori bertentangan dengan banyak
intuisi populer tentang hukum bilangan besar. Karena bahkan pembaca yang
berpengalaman dapat meragukan perilaku tersebut, ia menyertakan catatan
percobaan simulasi.\footnote{W.~Feller, \emx {Introduction to Probability Theory
and its Applications,} vol.~1, 3rd~ed. (New York: John Wiley \& Sons, 1968),
p.~xi. Ringkasan editorial; ungkapan sumber tidak diterjemahkan atau
direproduksi.}

\noindent Feller memberikan grafik hasil 10{,}000 permainan \emx
{kepala atau ekor} yang serupa dengan Gambar~\ref{fig 1.4.2}.

Sistem taruhan martingale\index{sistem taruhan martingale} 
yang dijelaskan dalam Latihan~\ref{exer 1.1.10} 
memiliki sejarah yang panjang dan menarik.  Russell Barnhart\index{BARNHART, R.} memberi tahu para
penulis bahwa penggunaannya setidaknya dapat dilacak hingga 1754, ketika Casanova\index{CASANOVA, G.}
menulis dalam memoarnya, \emx {History of My Life}. Ia menceritakan penggunaan
martingale di Ridotto di Venesia dengan terus menggandakan taruhan, memperoleh
tiga atau empat kemenangan per hari selama sisa Karnaval, dan tidak pernah
mencapai kekalahan pada taruhan keenam; menurut perhitungannya, kekalahan pada
tahap itu akan menghabiskan modal dua ribu zecchini.\footnote{G.~Casanova,
\emx {History of My Life,} vol.~IV, Chap.~7, trans.\ W.~R. Trask (New York:
Harcourt-Brace, 1968), p.~124. Ringkasan editorial; ungkapan terjemahan modern
tidak diterjemahkan ulang atau direproduksi.}

Bahkan seandainya tidak ada angka nol pada roda rolet sehingga permainan itu sepenuhnya
adil, sistem martingale, ataupun sistem lain apa pun, tidak dapat mengubahnya
menjadi permainan yang menguntungkan.  Gagasan bahwa permainan adil tetap adil dan
permainan tidak adil tetap tidak adil di bawah sistem perjudian telah dimanfaatkan oleh para matematikawan 
untuk memperoleh hasil-hasil penting dalam kajian peluang.  Kita akan memperkenalkan konsep umum 
martingale dalam Bab~\ref{chp 6}.

Kata \emx {martingale}\index{martingale!asal kata} sendiri juga memiliki sejarah yang 
menarik.  Asal-usul kata tersebut tidak jelas.  Edisi baru \emx {Oxford English Dictionary} memberikan contoh
penggunaannya pada awal 1600-an dan menyatakan bahwa kata itu mungkin berasal dari rujukan
dalam Buku Satu, Bab 20, karya Rabelais\index{RABELAIS, F.}. Dalam adegan itu,
Gargantua mempertimbangkan beberapa potongan celana bagi seorang orator,
termasuk potongan bernama \emx{martingale}, di samping gaya pelaut, Swiss, dan
``ekor ikan kod."\footnote{Dikutip dalam \emx {Portable Rabelais,}
ed.~S.~Putnam (New York: Viking, 1946), p.~113. Ringkasan editorial; ungkapan
terjemahan modern tidak diterjemahkan ulang atau direproduksi.}

Dominic Lusinchi\index{LUSINCHI, D.} mencatat kemunculan kata martingale yang lebih awal.  Menurut kamus bahasa Prancis \emx{Le Petit Robert},
kata tersebut berasal dari kata Proven\c{c}al ``martegalo," yang berarti ``dari Martigues."  Martigues adalah sebuah kota tepat di sebelah barat Merseille.  Kamus itu memberikan contoh ``chausses \`{a} la martinguale" (yang berarti celana bergaya Martigues) dan tahun 1491.  

Dalam penggunaan modern, martingale memiliki beberapa arti berbeda selain penggunaannya dalam perjudian; semuanya berkaitan dengan tindakan \emx
{menahan}.  Sebagai contoh, martingale adalah tali pada
perlengkapan kuda yang digunakan untuk menahan kepala kuda, dan juga bagian dari
perlengkapan layar yang digunakan untuk menahan tiang cucur.

Sistem Labouchere\index{sistem taruhan Labouchere} yang dijelaskan dalam Latihan~\ref{exer 1.1.9} 
dinamai menurut Henry du~Pre Labouchere\index{LABOUCHERE, H. du P.} 
(1831--1912), seorang wartawan Inggris dan anggota Parlemen.  La\-bou\-chere mengaitkan 
sistem tersebut dengan Condorcet\index{CONDORCET, Le Marquis de}.  Condorcet (1743--1794) adalah seorang pemimpin 
politik pada masa Revolusi Prancis yang tertarik menerapkan teori peluang 
pada ekonomi dan politik.  Sebagai contoh, ia menghitung peluang bahwa sebuah juri yang menggunakan
suara mayoritas akan memberikan keputusan yang benar apabila setiap juri memiliki peluang yang sama untuk
memutuskan dengan benar.  Tulisan-tulisannya menyediakan banyak gagasan tentang bagaimana
peluang dapat diterapkan pada urusan manusia.\footnote{ Le Marquise de~Condorcet, 
\emx {Essai sur l'Application de l'Analyse \`a la Probabilit\'e d\`es D\'ecisions
Rendues a la Pluralit\'e des Voix} (Paris: Imprimerie Royale, 1785).}

\exercises
\begin{LJSItem}

\i\label{exer 1.1.1} Ubah program {\bf CoinTosses} agar melempar sebuah koin sebanyak $n$
kali dan, setelah setiap 100 pelemparan, mencetak proporsi kepala dikurangi 1/2.  Apakah 
bilangan-bilangan ini tampak mendekati 0 ketika $n$ meningkat?  Ubah lagi program tersebut agar,
setiap 100 kali, mencetak kedua besaran berikut: proporsi kepala 
dikurangi 1/2, dan jumlah kepala dikurangi setengah jumlah pelemparan.  Apakah bilangan-bilangan ini 
tampak mendekati 0 ketika $n$ meningkat?

\i\label{exer 1.1.2} Ubah program {\bf CoinTosses} agar melempar sebuah
koin sebanyak $n$ kali
dan mencatat apakah proporsi kepala berjarak tidak lebih dari .1 terhadap .5 (yaitu,
berada antara .4 dan~.6).  Minta program Anda mengulangi percobaan ini 100 kali. 
Kira-kira
seberapa besar $n$ harus dipilih agar pada sekitar 95 dari 100 pengulangan,
proporsi
kepala berada antara .4 dan .6?

\i\label{exer 1.1.3} Pada awal 1600-an, Galileo\index{GALILEO, G.} diminta menjelaskan fakta
bahwa,
meskipun banyaknya tripel bilangan bulat dari 1 sampai 6 yang berjumlah 9 sama
dengan
banyaknya tripel semacam itu yang berjumlah 10, ketika tiga dadu dilempar, jumlah 9 tampaknya
muncul lebih jarang daripada jumlah 10---setidaknya menurut pengalaman para penjudi.

\begin{enumerate}
\item Tulis program untuk menyimulasikan pelemparan tiga dadu berkali-kali
dan mencatat proporsi pelemparan yang jumlahnya 9 serta
proporsi pelemparan yang jumlahnya 10.
\item Dapatkah Anda menyimpulkan dari simulasi tersebut bahwa para penjudi itu benar?
\end{enumerate}

\i\label{exer 1.1.4} Dalam racquetball\index{racquetball}, seorang pemain terus melakukan servis selama
ia menang; sebuah
poin hanya diperoleh ketika seorang pemain sedang melakukan servis dan memenangkan reli.  Pemain pertama
yang memperoleh 21 poin memenangkan permainan.  Asumsikan bahwa Anda melakukan servis lebih dahulu dan memiliki
peluang .6 untuk memenangkan reli ketika Anda melakukan servis dan peluang .5 ketika
lawan Anda
melakukan servis.  Perkirakan, melalui simulasi, peluang bahwa Anda akan memenangkan
sebuah
permainan.

\i\label{exer 1.1.5} Pertimbangkan taruhan bahwa ketiga dadu akan memperlihatkan
angka enam setidaknya sekali dalam $n$ kali pelemparan tiga dadu.  Hitung $f(n)$, yaitu peluang 
setidaknya satu tripel-enam ketika tiga dadu dilempar $n$ kali.  
Tentukan nilai terkecil $n$ yang diperlukan agar taruhan bahwa 
tripel-enam akan muncul ketika tiga dadu dilempar $n$ kali menjadi menguntungkan.  (DeMoivre akan
mengatakan bahwa nilainya
seharusnya sekitar $216\log 2 = 149.7$ dan karena itu akan menjawab 150---lihat
Latihan~\ref{sec 1.2}.\ref{exer 1.2.16}.
Apakah Anda sependapat dengannya?)                           

\i\label{exer 1.1.6} Di Las Vegas, sebuah roda rolet\index{rolet} memiliki 38 petak
bernomor 0, 00, 1, 2, \dots, 36.  Petak 0 dan 00 berwarna hijau, sedangkan separuh dari
36 petak
yang tersisa berwarna merah dan separuhnya berwarna hitam.  Seorang bandar memutar roda
dan
melemparkan bola gading ke dalamnya.  Jika Anda mempertaruhkan 1 dolar pada merah, Anda memenangkan 1 dolar jika
bola berhenti di petak merah; jika tidak, Anda kehilangan 1 dolar.  Tulis program untuk
menentukan total perolehan seorang pemain yang memasang 1000 taruhan pada merah.

\i\label{exer 1.1.7} Bentuk taruhan lain dalam rolet adalah bertaruh bahwa suatu
angka tertentu (misalnya 17) akan muncul.  Jika bola berhenti pada angka Anda, Anda
mendapatkan kembali satu dolar Anda ditambah 35 dolar.  Jika tidak, Anda kehilangan satu dolar tersebut.  Tulis
program yang menggambarkan perolehan Anda ketika bermain rolet 500 kali di
Las
Vegas: pertama, ketika Anda setiap kali bertaruh pada merah (lihat Latihan~\ref{exer 1.1.6}), 
kemudian, pada kunjungan kedua ke Las Vegas, ketika Anda bermain 500 kali dan setiap
kali
bertaruh pada angka 17.  Perbedaan apa yang Anda lihat pada grafik perolehan
dalam
kedua kesempatan tersebut?

\i\label{exer 1.1.8} Seorang mahasiswa yang cermat memperhatikan bahwa, dalam simulasi kita atas
permainan kepala
atau ekor (lihat Contoh~\ref{exam 1.3}), proporsi permainan ketika
pemain selalu unggul sangat dekat dengan proporsi permainan ketika
perolehan total
pemain berakhir pada~0.  Hitung peluang-peluang ini dengan
mencacah
semua kasus untuk dua pelemparan dan empat pelemparan, lalu tentukan apakah menurut Anda
kedua peluang tersebut memang sama.

\i\label{exer 1.1.9} \emx {Sistem Labouchere}\index{sistem taruhan Labouchere} 
untuk rolet dimainkan
sebagai berikut.  Tuliskan sebuah daftar bilangan, biasanya 1, 2, 3, 4.  Pertaruhkan jumlah
bilangan
pertama dan terakhir, $1 + 4 = 5$, pada merah.  Jika menang, hapus bilangan pertama dan
terakhir
dari daftar Anda.  Jika kalah, tambahkan besar taruhan terakhir ke
akhir daftar.  Kemudian gunakan daftar baru itu dan pertaruhkan jumlah bilangan pertama dan
terakhir
(jika hanya ada satu bilangan, pertaruhkan sebesar bilangan tersebut).  Lanjutkan sampai
daftar Anda kosong.  Tunjukkan bahwa, jika hal ini terjadi, Anda memenangkan jumlah, $1 + 2 + 3
+ 4
= 10$, dari daftar awal Anda.  Simulasikan sistem ini dan periksa apakah Anda memang selalu
berhenti dan, dengan demikian, selalu menang.  Jika ya, mengapa sistem ini bukan sistem perjudian yang
selalu berhasil?

\i\label{exer 1.1.10} Sistem perjudian terkenal lainnya adalah \emx
{sistem penggandaan martingale}\index{sistem taruhan martingale}.  
Misalkan Anda bertaruh bahwa merah akan muncul
dalam rolet.  Setiap kali menang, pertaruhkan 1 dolar pada permainan berikutnya.  Setiap kali
kalah,
gandakan taruhan sebelumnya.  Misalkan Anda menggunakan sistem ini sampai telah memenangkan sedikitnya 5
dolar atau telah kehilangan lebih dari 100 dolar.  Tulis program untuk menyimulasikan
sistem ini, mainkan beberapa kali, lalu lihat hasilnya.  Dalam bukunya
\emx
{The Newcomes,} W.~M.~Thackeray\index{THACKERAY, W. M.} berkomentar, 
``Anda belum pernah bermain?  Jangan
bermain; terutama, hindarilah martingale jika Anda bermain."\footnote{W.~M.~Thackerey, \emx
{The
Newcomes} (London: Bradbury and Evans, 1854--55).}  Apakah ini nasihat yang baik?

\i\label{exer 1.1.11} Ubah program {\bf HTSimulation} agar
mencatat
perolehan maksimum Peter dalam setiap permainan yang terdiri atas 40 pelemparan.  Minta program Anda
mencetak
proporsi permainan ketika perolehan total bernilai $0,\ 2,\ 
4,\ \dots,\ 40$.  Hitung peluang eksak yang bersesuaian untuk permainan dengan
dua
pelemparan dan empat pelemparan.

\i\label{exer 1.1.12} 
Dalam pemilihan nasional mendatang untuk Presiden Amerika Serikat, seorang
peneliti jajak pendapat berencana memprediksi pemenang suara rakyat dengan mengambil
sampel acak
sebanyak 1000 pemilih dan menyatakan bahwa pemenangnya adalah kandidat yang memperoleh
suara terbanyak
dalam sampelnya. Misalkan 48 persen pemilih berencana memilih
kandidat
Partai Republik dan 52 persen berencana memilih kandidat Partai
Demokrat.  
Untuk memperoleh gambaran tentang kewajaran rencana tersebut, tulis program yang membuat
prediksi ini
melalui simulasi.  Ulangi simulasi 100 kali dan lihat berapa kali prediksi peneliti tersebut
terbukti benar.  Ulangi percobaan Anda, kini dengan mengasumsikan bahwa 49 persen
populasi berencana memilih kandidat Partai Republik; pertama dengan sampel 1000, kemudian
dengan sampel 3000.  (Gallup\index{Gallup Poll} Poll menggunakan sekitar 3000.)  (Gagasan ini
dibahas lebih lanjut dalam Bab~\ref{chp 9}, Bagian~\ref{sec 9.1}.)

\i\label{exer 1.1.13} Tversky\index{TVERSKY, A.} dan para koleganya melaporkan
bahwa sekitar empat dari lima responden memilih jawaban (a) pada persoalan ukuran
sampel berikut.\footnote{Rujukan historis: K.~McKean, ``Decisions, Decisions,"
\emx {Discover,} June 1985, pp.~22--31, yang melaporkan karya Tversky et.\ al.\
dalam \emx {Judgement Under Uncertainty: Heuristics and Biases} (Cambridge:
Cambridge University Press, 1982). Tidak ada teks artikel tersebut yang
diterjemahkan atau direproduksi dalam rumusan latihan ini.}

Sebuah kota dilayani dua rumah sakit\index{rumah sakit}. Rumah sakit besar mencatat
sekitar 45 kelahiran per hari dan rumah sakit kecil sekitar 15. Secara keseluruhan,
proporsi bayi laki-laki sekitar 50 persen, tetapi proporsi harian dapat menyimpang.
Selama satu tahun, rumah sakit mana yang diperkirakan mencatat lebih banyak hari ketika
lebih dari 60 persen bayi yang lahir adalah laki-laki?
\begin{enumerate}
\item rumah sakit besar
\item rumah sakit kecil
\item tidak keduanya---jumlah harinya akan kurang lebih sama.
\end{enumerate}

Asumsikan peluang kelahiran bayi laki-laki adalah~.5; taksiran sebenarnya lebih dekat
ke~.513. Tentukan jawabannya melalui simulasi. Dapatkah Anda menjelaskan hasilnya
dan mengapa banyak orang keliru dengan membandingkan variasi proporsi sampel
untuk ukuran sampel 45 dan 15?

\i\label{exer 1.1.14} Anda ditawari permainan berikut.  Sebuah koin seimbang akan
dilempar sampai
sisi kepala muncul untuk pertama kalinya.  Jika hal ini terjadi pada pelemparan ke-$j$, Anda dibayar
$2^j$ dolar.  Anda pasti memenangkan setidaknya 2 dolar, jadi semestinya Anda
bersedia
membayar untuk memainkan permainan ini---tetapi berapa banyak?  Hanya sedikit orang yang bersedia membayar sebanyak 10
dolar untuk memainkannya.  Cobalah menentukan, melalui simulasi, jumlah
wajar
yang bersedia Anda bayar per permainan jika Anda diperbolehkan
memainkannya berkali-kali.
Apakah jumlah yang bersedia Anda bayar per permainan bergantung pada
banyaknya permainan yang diperbolehkan?

\i\label{exer 1.1.15} Tversky dan para koleganya menelaah catatan 48 pertandingan
Philadelphia 76ers\index{Philadelphia 76ers} pada musim 1980--81. Para pemain
memperkirakan bahwa peluang tembakan berhasil setelah keberhasilan sebelumnya sekitar
25 persen lebih besar daripada setelah kegagalan; catatan pertandingan justru menunjukkan
peluang berhasil 6 persen lebih besar setelah kegagalan. Banyaknya rentetan panas dan dingin
dinilai sesuai dengan variasi acak.\footnote{Ringkasan editorial berdasarkan K.~McKean,
``Decisions, Decisions," \emx{Discover}, June 1985, pp.~22--31; ungkapan artikel tidak
diterjemahkan atau direproduksi.} Rumuskan model sederhana: dalam setiap pertandingan,
seorang pemain melakukan 20 tembakan yang saling bebas, dan setiap tembakan berhasil dengan
peluang~.5. Perkirakan melalui simulasi proporsi pertandingan yang memuat sedikitnya satu
rentetan lima keberhasilan berturut-turut.

\i\label{1.1.16} Perkirakan, melalui simulasi, rata-rata jumlah anak
dalam suatu keluarga 
jika semua orang memiliki anak sampai memperoleh seorang anak laki-laki.  Lakukan hal yang sama
jika semua
orang memiliki anak sampai memperoleh setidaknya satu anak laki-laki dan setidaknya satu anak perempuan. 
Berapa
banyak tambahan anak yang Anda perkirakan akan terdapat dalam skema kedua dibandingkan
dengan skema pertama pada 100{,}000 keluarga?  (Asumsikan bahwa anak laki-laki dan perempuan
memiliki peluang yang
sama.)

\putfig{4truein}{PSfig1-6}{Jalan acak.}{fig 1.5}

\i\label{exer 1.1.17} Pada awal 1900-an, George
P\'olya\index{P\'OLYA, G.} sering berjalan di hutan dekat tempat tinggalnya di
Zurich sambil memikirkan matematika. Dalam sebuah tulisan kemudian, ia
menceritakan secara garis besar bahwa pertemuan berulang yang tidak disengaja
dengan pasangan yang sama dalam satu perjalanan membuatnya memikirkan apakah
dua pejalan acak pasti akan bertemu.\footnote{Ringkasan editorial berdasarkan
G.~P\'olya, ``Two Incidents," \emx {Scientists at Work: Festschrift in Honour
of Herman Wold,} ed. T. Dalenius, G. Karlsson, dan S. Malmquist (Uppsala:
Almquist \& Wiksells Boktryckeri AB, 1970). Kutipan asli tidak direproduksi
atau diterjemahkan di sini.}
\par
P\'olya membahas pejalan acak dalam satu, dua, dan tiga dimensi\index{jalan acak!dalam $n$ 
dimensi}.  Dalam satu dimensi, ia membayangkan pejalan itu berada di sebuah jalan yang sangat panjang.  Pada setiap
persimpangan, pejalan tersebut melempar koin seimbang untuk menentukan arah langkah
berikutnya (lihat Gambar~\ref {fig 1.5}a).  Dalam dua dimensi, pejalan tersebut berjalan
pada kisi jalan, dan di setiap persimpangan ia memilih salah satu dari empat
arah yang mungkin dengan peluang sama (lihat Gambar~\ref {fig 1.5}b).  Dalam
tiga dimensi (mungkin lebih tepat jika kita menyebutnya pemanjat acak), pejalan tersebut
bergerak pada kisi tiga dimensi, dan di setiap persimpangan kini terdapat enam
arah berbeda yang dapat dipilih, masing-masing dengan peluang sama
(lihat Gambar~\ref {fig 1.5}c).
\par
Pembaca dipersilakan melihat Bagian~\ref{sec 12.1}, tempat masalah ini dan masalah-masalah terkait 
dibahas.

\begin{enumerate}

\item Tulis program untuk menyimulasikan jalan acak satu dimensi yang dimulai dari
0.  Minta program mencetak panjang selang waktu antara setiap kembalinya pejalan ke
titik
awal (kembali ke 0).  Periksa apakah dari simulasi ini Anda dapat menebak
jawaban atas pertanyaan berikut: Apakah pejalan pada akhirnya selalu kembali ke
titik
awal, atau mungkinkah ia terus menjauh selamanya?

\item Lintasan dua pejalan dalam dua dimensi yang bertemu setelah $n$ langkah
dapat dipandang sebagai satu lintasan yang dimulai dari $(0,0)$ dan kembali ke
$(0,0)$ setelah $2n$ langkah.  Ini berarti bahwa peluang dua pejalan
acak dalam dua dimensi bertemu sama dengan peluang bahwa satu
pejalan dalam dua dimensi pernah kembali ke titik awal.  Jadi,
pertanyaan apakah dua pejalan pasti bertemu sama dengan pertanyaan
apakah
satu pejalan pasti kembali ke titik awal.

Tulis program untuk menyimulasikan jalan acak dalam dua dimensi dan tentukan apakah menurut Anda
pejalan itu pasti kembali ke $(0,0)$.  Jika ya, P\'olya tentu
akan
terus bertemu teman-temannya di taman.  Mungkin sekarang Anda telah membuat dugaan 
atas pertanyaan berikut: Apakah pejalan acak dalam satu atau dua dimensi pasti
kembali
ke titik awal?  P\'olya menjawab pertanyaan ini untuk dimensi
satu, dua, dan tiga.  Ia membuktikan hasil luar biasa bahwa jawabannya adalah
\emx {ya} dalam satu dan dua dimensi, serta \emx {tidak} dalam tiga dimensi.

\item Tulis program untuk menyimulasikan jalan acak dalam tiga dimensi dan tentukan
apakah, berdasarkan simulasi ini serta hasil pada (a) dan (b), 
Anda dapat menebak hasil P\'olya.  
\end{enumerate}

\end{LJSItem}

\section{Distribusi Peluang Diskret}\label{sec 1.2}  

Dalam buku ini kita akan mempelajari berbagai percobaan dari sudut pandang
peluang.  Cakupan kajian ini akan menjadi jelas seiring dengan dikembangkannya
teori dan dianalisisnya contoh-contoh.  Namun, gagasan umumnya dapat
dijelaskan dan digambarkan sebagai berikut: dengan setiap percobaan yang kita bahas
akan dikaitkan suatu peubah acak yang menyatakan hasil dari
percobaan tertentu.  Himpunan semua hasil yang mungkin disebut \emx
{ruang sampel}\index{ruang sampel}.  Pada bagian awal ini, kita akan membahas
kasus ketika percobaan hanya memiliki sejumlah terhingga hasil yang mungkin, yaitu ruang sampelnya
terhingga.  Kemudian kita akan memperumum ke kasus ketika ruang sampel terhingga atau
takhingga terhitung.  Hal ini membawa kita pada definisi berikut.

\subsection*{Peubah Acak dan Ruang Sampel}

\begin{definition}\label{def 1.1} 
Misalkan kita memiliki sebuah percobaan yang hasilnya bergantung pada faktor acak.
Kita menyatakan hasil percobaan tersebut dengan
huruf Romawi kapital, seperti $X$, yang disebut \emx {peubah acak}\index{peubah acak}.
\emx {Ruang sampel}\index{ruang sampel} percobaan adalah himpunan semua 
hasil yang mungkin.  Jika ruang sampel terhingga atau takhingga terhitung, peubah acak
tersebut disebut \emx {diskret}.\index{peubah acak!diskret}
\end{definition}
Umumnya kita menyatakan ruang sampel dengan huruf Yunani kapital $\Omega$.  Seperti
dinyatakan di atas, dalam korespondensi antara suatu percobaan dan
teori matematika yang digunakan untuk mengkajinya, ruang sampel $\Omega$ bersesuaian dengan
himpunan hasil yang mungkin dari percobaan tersebut.    
\par
Sekarang kita memberikan dua definisi tambahan.  Definisi ini melengkapi
definisi ruang sampel dan memperjelas beberapa istilah umum yang digunakan
bersama ruang sampel.  Pertama, elemen-elemen suatu ruang sampel kita sebut
\emx {hasil}\index{hasil}.  Kedua, setiap himpunan bagian dari ruang sampel disebut 
\emx {kejadian}\index{kejadian}. Biasanya, hasil akan kita nyatakan dengan huruf kecil dan
kejadian dengan huruf kapital.
 
\begin{example}\label{exam 1.5}
Sebuah dadu dilempar sekali.  Misalkan $X$ menyatakan hasil percobaan ini.  Maka ruang
sampel percobaan ini adalah himpunan dengan 6 elemen
$$
\Omega = \{1,2,3,4,5,6\}\ ,
$$
di mana setiap hasil $i$, untuk~$i = 1$, \dots, 6, bersesuaian dengan banyaknya
titik pada sisi yang muncul.  Kejadian
$$
E = \{2,4,6\}
$$
bersesuaian dengan pernyataan bahwa hasil pelemparan merupakan bilangan genap.  Kejadian $E$
juga dapat dijelaskan dengan mengatakan bahwa $X$ genap.  Kecuali ada alasan untuk meyakini bahwa dadu
tersebut tidak seimbang, asumsi yang wajar adalah bahwa setiap hasil memiliki peluang sama.  Menggunakan
kesepakatan ini berarti kita menetapkan peluang 1/6 pada masing-masing dari keenam hasil, yaitu $m(i) =
1/6$, untuk $1 \le i \le 6$.
\end{example}

\subsection*{Fungsi Distribusi}

Selanjutnya kita menjelaskan penetapan peluang.  Definisi-definisi berikut
dimotivasi oleh contoh di atas, ketika kita memasangkan pada setiap hasil dalam
ruang sampel sebuah bilangan taknegatif sedemikian rupa sehingga jumlah semua bilangan yang dipasangkan
sama dengan 1.

\begin{definition}\label{def 1.2} 
Misalkan $X$ adalah peubah acak yang menyatakan nilai hasil suatu percobaan tertentu,
dan asumsikan bahwa percobaan ini hanya memiliki sejumlah terhingga hasil yang mungkin.  Misalkan $\Omega$ adalah
ruang sampel percobaan (yaitu himpunan semua nilai yang mungkin dari $X$, atau secara ekuivalen,
himpunan semua hasil yang mungkin dari percobaan tersebut).  \emx {Fungsi
distribusi}\index{fungsi distribusi} bagi $X$ adalah fungsi bernilai
real $m$ yang berdomain $\Omega$ dan memenuhi:
\begin{enumerate}
\item $m(\omega) \geq 0\ , \qquad$untuk semua $\,\,\omega\in\Omega\ $, dan  
\medbreak
\item $\sum\limits_{\omega \in \Omega} m(\omega) = 1\ $.  
\end{enumerate}
Untuk setiap himpunan bagian $E$ dari $\Omega$, kita mendefinisikan \emx {peluang}\index{peluang!suatu kejadian}
dari
$E$ sebagai bilangan $P(E)$ yang diberikan oleh
$$
P(E) = \sum_{\omega\in E} m(\omega)\ .
$$
\end{definition}
 
\begin{example}\label{exam 1.6}
Pertimbangkan percobaan ketika sebuah koin dilempar dua kali.  Misalkan $X$ adalah peubah acak yang
bersesuaian dengan percobaan ini.  Perhatikan bahwa terdapat beberapa cara untuk mencatat hasil
percobaan tersebut.  Sebagai contoh, kita dapat mencatat kedua pelemparan menurut urutan
terjadinya.  Dalam hal ini, kita memiliki $\Omega = $\{HH,HT,TH,TT\}.  Kita juga dapat mencatat hasil
dengan hanya mencatat jumlah kepala yang muncul.  Dalam hal ini, kita memiliki $\Omega =
$\{0,1,2\}.  Terakhir, kita dapat mencatat kedua hasil tanpa memperhatikan urutan
terjadinya.  Dalam hal ini, kita memiliki $\Omega = $\{HH,HT,TT\}.
\par
Untuk sementara, kita akan menggunakan ruang sampel pertama yang diberikan di atas.  Kita akan mengasumsikan bahwa
keempat hasil memiliki peluang sama, dan mendefinisikan fungsi distribusi $m(\omega)$ dengan
$$
m(\mbox{HH}) = m(\mbox{HT}) = m(\mbox{TH}) = m(\mbox{TT}) = \frac14\ .
$$
\par
Misalkan $E = $\{HH,HT,TH\} adalah kejadian bahwa setidaknya satu kepala muncul.  Maka peluang
$E$ dapat dihitung sebagai berikut: 
\begin{eqnarray*}
P(E) &=& m(\mbox{HH}) + m(\mbox{HT}) + m(\mbox{TH}) \\
     &=& \frac14 + \frac14 + \frac14 = \frac34\ .
\end{eqnarray*}

Demikian pula, jika $F = $\{HH,HT\} adalah kejadian bahwa kepala muncul pada pelemparan pertama, maka kita memperoleh
\begin{eqnarray*}
P(F) &=& m(\mbox{HH}) + m(\mbox{HT}) \\
     &=& \frac14 + \frac14 = \frac12\ .
\end{eqnarray*}
\end{example}
 
\begin{example}\label{1.8}(Lanjutan Contoh \ref{exam 1.5})
Ruang sampel percobaan pelemparan dadu adalah
himpunan dengan 6 elemen
$\Omega = \{1,2,3,4,5,6\}$.  Kita mengasumsikan bahwa dadu itu seimbang dan memilih
fungsi distribusi yang didefinisikan oleh
$$
m(i) = \frac16, \qquad {\rm{untuk}}\,\, i = 1, \dots, 6\ .
$$
Jika $E$ adalah kejadian bahwa hasil pelemparan merupakan bilangan genap, maka $E =
\{2,4,6\}$ dan
\begin{eqnarray*}
P(E) &=& m(2) + m(4) + m(6) \\
     &=& \frac16 + \frac16 + \frac16 = \frac12\ .
\end{eqnarray*}
\end{example}
\par
Perhatikan bahwa definisi di atas langsung mengakibatkan bahwa, untuk
setiap $\omega \in \Omega$,
$$
P(\{\omega\}) = m(\omega)\ .
$$
Artinya, peluang kejadian elementer $\{\omega\}$, yang terdiri atas
satu hasil $\omega$, sama dengan nilai $m(\omega)$ yang dipasangkan pada
hasil $\omega$ oleh fungsi distribusi.
 
\begin{example}\label{exam 1.7}
Tiga orang, A,~B, dan~C, mencalonkan diri untuk jabatan yang sama, dan kita mengasumsikan
bahwa
tepat satu dari mereka menang.  Ruang sampelnya dapat dipilih sebagai
himpunan dengan 3 elemen
$\Omega = $\{A,B,C\}, dengan setiap elemen bersesuaian dengan hasil berupa kemenangan
kandidat tersebut.  Misalkan A dan B memiliki peluang menang yang sama,
tetapi peluang C hanya 1/2 dari peluang A atau B.  Maka kita menetapkan
$$
m(\mbox{A}) = m(\mbox{B}) = 2m(\mbox{C})\ .
$$
Karena
$$
m(\mbox{A}) + m(\mbox{B}) + m(\mbox{C}) = 1\ ,
$$
kita melihat bahwa
$$
2m(\mbox{C}) + 2m(\mbox{C}) + m(\mbox{C}) = 1\ ,
$$
yang menyiratkan bahwa $5m(\mbox{C}) = 1$.  Oleh karena itu,
$$
m(\mbox{A}) = \frac25\ , \qquad m(\mbox{B}) = \frac25\ , \qquad m(\mbox{C}) =
\frac15\ .
$$
Misalkan $E$ adalah kejadian bahwa A atau C menang.  Maka $E = $\{A,C\}, dan
$$
P(E) = m(\mbox{A}) + m(\mbox{C}) = \frac25 + \frac15 = \frac35\ .
$$
\end{example}
\par
Dalam banyak kasus, kejadian dapat dijelaskan melalui kejadian lain dengan menggunakan
konstruksi baku teori himpunan.  Kita akan meninjau secara singkat definisi
konstruksi-konstruksi tersebut.  Lihat Gambar~\ref{fig 1.6} untuk diagram Venn
yang menggambarkan konstruksi ini. 
\par
Misalkan $A$ dan $B$ adalah dua himpunan.  Gabungan $A$ dan $B$ adalah himpunan
$$A \cup B = \{x\,|\, x \in A\ \mbox{atau}\ x \in B\}\ .$$
Irisan $A$ dan $B$ adalah himpunan
$$A \cap B = \{x\,|\, x \in A\ \mbox{dan}\ x \in B\}\ .$$
Selisih $A$ dan $B$ adalah himpunan
$$A - B = \{x\,|\, x \in A\ \mbox{dan}\ x \not \in B\}\ .$$
Himpunan $A$ adalah himpunan bagian dari $B$, ditulis 
$A \subset B$, jika setiap elemen $A$ juga merupakan elemen $B$.  Terakhir,
komplemen $A$ adalah himpunan
$$\tilde A = \{x\,|\, x \in \Omega\ \mbox{dan}\ x \not \in A\}\ .$$
\par
Konstruksi ini penting karena biasanya
kejadian rumit yang dijelaskan dengan bahasa sehari-hari dapat diuraikan menjadi kejadian yang lebih sederhana
dengan menggunakan konstruksi tersebut.  Sebagai contoh, jika $A$ adalah kejadian ``besok akan turun salju
dan lusa akan
turun hujan," $B$ adalah kejadian ``besok akan turun salju," dan
$C$ adalah
kejadian ``dua hari lagi akan turun hujan," maka $A$ merupakan irisan
kejadian
$B$ dan $C$.  Demikian pula, jika $D$ adalah kejadian ``besok akan turun salju atau
lusa akan turun hujan," maka $D = B \cup C$.  (Perhatikan bahwa kita harus
berhati-hati di sini,
sebab kadang-kadang kata ``atau" berarti tepat satu dari kedua
alternatif
akan terjadi.  Maknanya biasanya jelas dari konteks.  Dalam buku ini, kita
selalu menggunakan
kata ``atau" dalam arti inklusif, yaitu $A$ atau $B$ berarti setidaknya
satu dari kedua
kejadian $A$, $B$ terjadi.)  Kejadian $\tilde B$ adalah kejadian ``besok
tidak turun
salju."  Terakhir, jika $E$ adalah kejadian ``besok akan turun salju, tetapi lusa
tidak akan turun hujan," maka $E = B - C$. 

\putfig{5truein}{PSfig1-7}{Operasi dasar himpunan.}{fig 1.6} 


\subsection*{Sifat-sifat}
\nopagebreak
\begin{theorem}\label{thm 1.1}
Peluang yang dipasangkan pada kejadian oleh suatu fungsi distribusi\index{fungsi distribusi!
sifat-sifat} pada ruang sampel
$\Omega$ 
memenuhi sifat-sifat berikut: 
\begin{enumerate}
\item  $P(E) \geq 0$ untuk setiap $E \subset \Omega\ $.
\item  $P(\Omega) = 1\ $.
\item  Jika $E \subset F \subset \Omega$, maka $P(E) \leq P(F)\ $.
\item  Jika $A$ dan $B$ adalah himpunan bagian \emx {saling lepas} dari $\Omega$, maka $P(A
\cup B)
= P(A) + P(B)\ $.
\item  $P(\tilde A) = 1 - P(A)$ untuk setiap $A \subset \Omega\ $.
\end{enumerate}
\proof
Untuk setiap kejadian $E$, peluang $P(E)$ ditentukan dari distribusi $m$
oleh
$$
P(E) = \sum_{\omega \in E} m(\omega)\ ,
$$
untuk setiap $E \subset \Omega$.
Karena fungsi $m$ taknegatif, $P(E)$ juga
taknegatif.  Jadi, Sifat 1 benar.  
\par
Sifat 2 dibuktikan oleh persamaan
$$
P(\Omega) = \sum_{\omega \in \Omega} m(\omega) = 1\ .
$$
\par
Misalkan $E \subset F \subset \Omega$.  Maka setiap elemen $\omega$ yang
termasuk dalam $E$ juga termasuk dalam $F$.  Oleh karena itu,
$$
\sum_{\omega \in E} m(\omega) \leq \sum_{\omega \in F} m(\omega)\ ,
$$
karena setiap suku dalam jumlah di sebelah kiri terdapat dalam jumlah di sebelah kanan, dan semua suku 
dalam kedua jumlah tersebut taknegatif.  Hal ini menyiratkan bahwa
$$P(E) \le P(F)\ ,$$
dan Sifat 3 terbukti.  
\par
Selanjutnya, misalkan $A$ dan $B$ adalah himpunan bagian yang saling lepas dari
$\Omega$.  Maka setiap elemen $\omega$ dari $A \cup B$ berada dalam $A$ dan
tidak berada
dalam $B$, atau berada dalam $B$ dan tidak berada dalam $A$.  Akibatnya,
$$\begin{array}{ll}
P(A \cup B) &= \sum_{\omega \in A \cup B} m(\omega) = \sum_{\omega \in A}
m(\omega) + \sum_{\omega \in B} m(\omega) \\
            &                             \\
            &= P(A) + P(B)\ ,
\end{array}
$$
dan Sifat 4 terbukti.  
\par
Terakhir, untuk membuktikan Sifat 5, perhatikan gabungan saling lepas
$$
\Omega = A \cup \tilde A\ .
$$
Karena $P(\Omega) = 1$, sifat aditivitas saling lepas (Sifat 4) menyiratkan
bahwa
$$
1 = P(A) + P(\tilde A)\ ,
$$
sehingga $P(\tilde A) = 1 - P(A)$.
\end{theorem}
\par
Penting untuk dipahami bahwa Sifat~4 dalam Teorema~\ref{thm 1.1} 
dapat diperluas untuk lebih dari dua himpunan.  Sifat aditivitas hingga
umum\index{sifat aditivitas hingga} diberikan oleh
teorema berikut.\hfil\break
\begin{theorem}\label{thm 1.1.5}
Jika $A_1$, \dots, $A_n$ merupakan himpunan bagian $\Omega$ yang saling lepas berpasangan
(yaitu, tidak ada dua dari himpunan $A_i$ yang memiliki elemen yang sama), maka
$$
P(A_1 \cup \cdots \cup A_n) = \sum_{i = 1}^n P(A_i)\ .
$$
\proof
Misalkan $\omega$ adalah sebarang elemen dalam gabungan
$$A_1 \cup \cdots \cup A_n\ .$$
Maka $m(\omega)$ muncul tepat satu kali pada masing-masing ruas kesamaan dalam
pernyataan teorema.
\end{theorem} 
\par
Kita akan sering menggunakan akibat berikut dari teorema di atas.
\begin{theorem}\label{thm 1.2}
 Misalkan $A_1$, \dots, $A_n$ adalah kejadian yang saling lepas berpasangan dengan $\Omega = A_1
\cup \cdots \cup A_n$, dan misalkan $E$ adalah sebarang kejadian.  Maka
$$
P(E) = \sum_{i = 1}^n P(E \cap A_i)\ .
$$
\proof
Himpunan $E \cap A_1$, \dots, $E \cap A_n$ saling lepas berpasangan, dan 
gabungannya adalah himpunan $E$.  Hasil tersebut kemudian mengikuti dari Teorema~\ref{thm 1.1.5}.
\end{theorem}
\begin{corollary}\label{cor 1.1}
Untuk sebarang dua kejadian $A$ dan $B$,
$$
P(A) = P(A \cap B) + P(A \cap \tilde B)\ .
$$
\end{corollary}
\par
Sifat~4 dapat diperumum dengan cara lain.  Misalkan $A$ dan $B$ adalah
himpunan bagian $\Omega$ yang tidak harus saling lepas.  Maka:\hfil\break
\begin{theorem}\label{thm 1.2.5}
Jika $A$ dan $B$ merupakan himpunan bagian dari $\Omega$, maka
\begin{equation}
P(A \cup B) = P(A) + P(B) - P(A \cap B)\ .
\label{eq 1.1}
\end{equation}                                                        
\proof
Ruas kiri Persamaan~\ref{eq 1.1} adalah jumlah $m(\omega)$ untuk     
$\omega$ yang berada dalam $A$ atau $B$.  Kita harus menunjukkan bahwa ruas kanan 
Persamaan~\ref{eq 1.1} juga menjumlahkan $m(\omega)$ untuk $\omega$ dalam $A$ atau $B$. 
Jika $\omega$ berada tepat dalam salah satu dari kedua himpunan, maka ia hanya dihitung dalam satu
dari tiga suku pada ruas kanan Persamaan~\ref{eq 1.1}.                
Jika ia berada dalam $A$ sekaligus $B$, ia dijumlahkan dua kali dari perhitungan
$P(A)$ dan $P(B)$
dan dikurangi satu kali untuk $P(A \cap B)$.  Jadi, ia dihitung tepat satu kali oleh
ruas kanan.  Tentu saja, jika $A \cap B = \emptyset$, Persamaan~\ref{eq 1.1} 
tereduksi menjadi Sifat~4.
(Persamaan~\ref{eq 1.1} juga dapat diperumum; lihat Teorema~\ref{thm 3.10}.)               
\end{theorem}


\subsection*{Diagram Pohon}

\begin{example}\label{exam 1.8}
Mari kita gambarkan sifat-sifat peluang kejadian melalui tiga kali
pelemparan sebuah koin.  Jika suatu percobaan berlangsung secara bertahap
seperti
ini, sering kali lebih mudah untuk menyatakan hasilnya dengan \emx {diagram
pohon}\index{diagram pohon} seperti ditunjukkan pada Gambar~\ref {fig 1.7}.

\putfig{4truein}{PSfig1-8}{Diagram pohon untuk tiga kali pelemparan koin.}{fig 1.7} 

Suatu \emx {lintasan} pada pohon bersesuaian dengan salah satu hasil yang mungkin dari
percobaan.  Untuk tiga kali pelemparan koin, terdapat delapan lintasan,
$\omega_1$,~$\omega_2$, \dots,~$\omega_8$; dengan mengasumsikan bahwa setiap hasil
berpeluang sama, kita menetapkan bobot yang sama, 1/8, pada setiap lintasan.  Misalkan $E$ adalah
kejadian
``setidaknya satu kepala muncul."  Maka $\tilde E$ adalah kejadian ``tidak ada kepala yang
muncul."  Kejadian ini hanya terjadi untuk satu hasil, yaitu $\omega_8 = \mbox{TTT}$. 
Jadi, $\tilde E = \{\mbox{TTT}\}$ dan kita memperoleh
$$
P(\tilde E) = P(\{\mbox{TTT}\}) = m(\mbox{TTT}) = \frac18\ .
$$
Menurut Sifat~5 dari Teorema~\ref{thm 1.1},
$$
P(E) = 1 - P(\tilde E) = 1 - \frac18 = \frac78\ .
$$
Perhatikan bahwa sering kali lebih mudah menghitung peluang bahwa suatu
kejadian tidak terjadi daripada peluang bahwa kejadian itu terjadi.  Selanjutnya kita menggunakan
Sifat~5 untuk memperoleh peluang yang diinginkan.

Misalkan $A$ adalah kejadian ``hasil pertama adalah kepala," dan $B$ adalah kejadian ``hasil
kedua adalah ekor."  Dengan melihat lintasan pada Gambar~\ref{fig 1.7},
kita memperoleh
$$
P(A) = P(B) = \frac12\ .
$$
Selain itu, $A \cap B = \{\omega_3,\omega_4\}$, sehingga $P(A \cap B) = 1/4.$ 
Dengan menggunakan Teorema~\ref{thm 1.2.5}, kita memperoleh
\begin{eqnarray*}
P(A \cup B) & = & P(A) + P(B) - P(A \cap B) \\
& = & \frac 12 + \frac 12 - \frac 14 = \frac 34\ .
\end{eqnarray*}
Karena $A \cup B$ merupakan himpunan dengan 6 elemen,
$$
A \cup B = \{\mbox{HHH,HHT,HTH,HTT,TTH,TTT}\}\ ,
$$
kita melihat bahwa pencacahan langsung memberikan hasil yang sama.
\end{example}
\par
Dalam contoh pelemparan koin dan contoh pelemparan dadu, kita telah menetapkan
peluang yang sama pada setiap hasil percobaan yang mungkin. 
Bersesuaian
dengan cara penetapan peluang ini, kita memiliki definisi berikut.

\subsection*{Distribusi Seragam}

\begin{definition}\label{def 1.3} 
\emx {Distribusi seragam}\index{distribusi seragam} pada ruang sampel $\Omega$ 
yang memuat $n$ elemen adalah fungsi $m$ yang didefinisikan oleh
$$
m(\omega) = \frac1n\ ,
$$
untuk setiap $\omega \in \Omega$.
\end{definition}
Penting untuk dipahami bahwa ketika suatu percobaan dianalisis untuk menjelaskan
hasil-hasil yang
mungkin, tidak ada satu pilihan ruang sampel yang selalu benar.  Untuk
percobaan
pelemparan koin dua kali dalam Contoh~\ref{exam 1.1},
kita memilih himpunan dengan 4 elemen $\Omega = $\{HH,HT,TH,TT\} sebagai ruang sampel
dan
menetapkan fungsi distribusi seragam.  Pilihan-pilihan ini tentu
wajar secara intuitif.  Namun, untuk tujuan tertentu mungkin lebih
berguna
mempertimbangkan ruang sampel dengan 3 elemen $\bar\Omega = \{0,1,2\}$, dengan
0
merupakan hasil ``tidak ada kepala yang muncul," 1 merupakan hasil ``tepat satu kepala
muncul," dan 2 merupakan hasil ``dua kepala muncul."  Fungsi distribusi
$\bar m$ pada $\bar\Omega$ yang didefinisikan oleh persamaan
$$
\bar m(0) = \frac14\ ,\qquad \bar m(1) = \frac12\ , \qquad \bar
m(2) =
\frac14
$$
adalah fungsi yang bersesuaian dengan kerapatan peluang seragam pada
ruang sampel asal $\Omega$.  Perhatikan bahwa kita sepenuhnya dapat memilih
fungsi distribusi lain.  Sebagai contoh, kita dapat mempertimbangkan fungsi
distribusi seragam pada $\bar\Omega$, yaitu fungsi $\bar q$
yang didefinisikan oleh
$$
\bar q(0) = \bar q(1) = \bar q(2) = \frac13\ .
$$
Meskipun $\bar q$ merupakan fungsi distribusi yang sah, fungsi ini tidak
sesuai dengan data pengamatan pelemparan koin.
\begin{example}\label{exam 1.11}
Pertimbangkan percobaan yang terdiri atas pelemparan sepasang dadu.  Sebagai
ruang
sampel $\Omega$, kita ambil himpunan semua pasangan terurut $(i,j)$ dari
bilangan bulat dengan $1\leq i\leq 6$ dan $1\leq j\leq 6$.  Jadi,
$$
\Omega = \{\,(i,j):1\leq i,\space j \leq 6\,\}\ .
$$
(Setidaknya terdapat satu pilihan ruang sampel lain yang ``wajar," yaitu
himpunan semua pasangan takterurut bilangan bulat, masing-masing antara 1 dan 6.  Untuk pembahasan
mengapa himpunan ini tidak digunakan, lihat Contoh~\ref{exam 3.15}.)
Untuk menentukan ukuran $\Omega$, perhatikan bahwa terdapat enam pilihan bagi $i$,
dan untuk setiap pilihan $i$ terdapat enam pilihan bagi $j$, sehingga menghasilkan 36
hasil berbeda.  Asumsikan bahwa dadu-dadu tersebut seimbang.  Dalam istilah matematika, hal ini
berarti
kita mengasumsikan setiap dari 36 hasil berpeluang sama, atau
secara ekuivalen,
kita menggunakan fungsi distribusi seragam pada $\Omega$ dengan menetapkan
$$
m((i,j)) = \frac1{36},\qquad 1\leq i,\space j \leq 6\ .
$$
Berapakah peluang memperoleh jumlah 7 dalam pelemparan dua dadu---atau
memperoleh jumlah 11?  Kejadian pertama, yang dinyatakan dengan $E$, adalah himpunan bagian
$$
E = \{(1,6),(6,1),(2,5),(5,2),(3,4),(4,3)\}\ .
$$
Jumlah 11 adalah himpunan bagian $F$ yang diberikan oleh
$$
F = \{(5,6),(6,5)\}\ .
$$
Akibatnya,
$$\begin{array}{ll}
P(E) = &\sum_{\omega \in E} m(\omega) = 6\cdot\frac1{36} = \frac16\ , \\
       &                                                            \\
P(F) = &\sum_{\omega \in F} m(\omega) = 2\cdot\frac1{36} = \frac1{18}\ .
\end{array}
$$

Berapakah peluang untuk tidak memperoleh \emx {mata ular (snakeeyes)}\index{mata ular} (dua angka satu)
maupun
\emx {sepasang enam (boxcars)}\index{sepasang enam} (dua angka enam)?  Kejadian memperoleh salah satu dari kedua
hasil tersebut adalah himpunan
$$
E = \{(1,1),(6,6)\}\ .
$$
Jadi, peluang untuk tidak memperoleh keduanya diberikan oleh
$$
P(\tilde E) = 1 - P(E) = 1 - \frac2{36} = \frac{17}{18}\ .
$$
\end{example}
\par
Dalam percobaan pelemparan koin dan dadu di atas, kita telah menetapkan
peluang yang sama pada setiap hasil.  Artinya, dalam setiap contoh kita telah memilih
fungsi distribusi seragam.  Pilihan ini wajar asalkan
koinnya seimbang dan dadu-dadunya tidak diberi pemberat.  Namun, keputusan mengenai
fungsi distribusi mana yang dipilih untuk menjelaskan suatu percobaan \emx {bukan} bagian
dari teori matematika dasar peluang.  Teori tersebut baru dimulai ketika
ruang sampel dan fungsi distribusi telah didefinisikan.

\subsection*{Penentuan Peluang}

Penting untuk mempertimbangkan bagaimana distribusi peluang
ditentukan dalam praktik.  Salah satu caranya adalah melalui \emx {simetri}.  Dalam kasus
pelemparan koin, kita tidak melihat perbedaan fisik antara kedua sisi
koin
yang semestinya memengaruhi peluang munculnya salah satu sisi. 
Demikian pula, pada dadu biasa tidak ada perbedaan hakiki antara
dua sisi mana pun, sehingga berdasarkan simetri kita menetapkan peluang yang sama bagi
setiap
hasil yang mungkin.  Secara umum, pertimbangan simetri sering menyarankan penggunaan
fungsi distribusi seragam.  Namun, kita harus berhati-hati.  Kita tidak selalu boleh mengasumsikan
bahwa,
hanya karena tidak mengetahui alasan yang menunjukkan bahwa satu hasil lebih
mungkin daripada yang lain, kita patut menetapkan peluang yang sama.  Sebagai
contoh, pertimbangkan percobaan menebak jenis kelamin bayi yang baru lahir.  Telah
diamati bahwa proporsi bayi baru lahir yang berjenis kelamin laki-laki kira-kira
.513.  Jadi, lebih tepat menggunakan fungsi distribusi yang
menetapkan peluang .513 pada hasil \emx {laki-laki} dan peluang .487 pada
hasil
\emx {perempuan} daripada menetapkan peluang 1/2 pada masing-masing hasil.  Ini merupakan
contoh ketika kita menggunakan pengamatan statistik untuk menentukan peluang. 
Perhatikan
bahwa peluang tersebut dapat berubah seiring adanya penelitian baru dan dapat berbeda dari
satu negara
ke negara lain.  Rekayasa genetika bahkan mungkin memungkinkan seseorang memengaruhi
peluang ini dalam kasus tertentu.

\subsection*{Rasio Peluang}\index{rasio peluang}

Perkiraan statistik untuk peluang dapat digunakan jika percobaan yang
dibahas dapat diulangi berkali-kali dalam keadaan serupa. 
Namun, misalkan pada awal musim sepak bola Amerika Anda ingin
menetapkan
peluang bagi kejadian bahwa Dartmouth\index{Dartmouth} akan mengalahkan Harvard\index{Harvard}.  
Anda sebenarnya tidak
memiliki data yang berkaitan dengan tim sepak bola tahun ini.  Meskipun demikian, Anda dapat
menentukan peluang pribadi dengan melihat jenis taruhan yang
bersedia
Anda pasang.  Sebagai contoh, misalkan Anda bersedia memasang taruhan 1
dolar dengan rasio peluang 2 banding 1 bahwa Dartmouth akan menang.  Artinya, Anda bersedia
membayar
2 dolar jika Dartmouth kalah untuk memperoleh 1 dolar jika
Dartmouth
menang.  Ini berarti bahwa menurut Anda peluang yang tepat bagi kemenangan Dartmouth
adalah 2/3.

Mari kita telaah lebih cermat hubungan antara rasio peluang dan peluang.
Misalkan kita bertaruh dengan rasio $r$ banding $1$ bahwa kejadian $E$ terjadi. 
Ini berarti kita menganggap bahwa $r$ kali lebih mungkin $E$ terjadi
daripada $E$
tidak terjadi.  Secara umum, rasio peluang $r$ banding $s$ diartikan sama
dengan
$r/s$ banding 1; dengan kata lain, hanya perbandingan antara kedua bilangan yang
penting ketika menyatakan rasio peluang.  
\par
Jika $r$ kali lebih mungkin $E$ terjadi daripada $E$
tidak terjadi, maka peluang $E$ terjadi haruslah $r/(r+1)$, sebab kita
memiliki
$$P(E) = r\,P(\tilde E)$$
dan
$$P(E) + P(\tilde E) = 1\ .$$
Secara umum, pernyataan bahwa rasio peluangnya adalah $r$ banding $s$ untuk mendukung kejadian
$E$
ekuivalen dengan pernyataan bahwa 
\begin{eqnarray*}
P(E) & = & \frac{r/s}{(r/s) + 1}\\
& = & \frac {r}{r+s}\ .
\end{eqnarray*}
Jika kita menetapkan $P(E) = p$, persamaan di atas dapat dengan mudah diselesaikan untuk $r/s$
sebagai fungsi dari
$p$; kita memperoleh $r/s = p/(1-p)$.  Pembahasan di atas dirangkum dalam
definisi berikut. 

\begin{definition}\label{def 1.4} 
Jika $P(E) = p$, \emx {rasio peluang} yang mendukung terjadinya kejadian $E$ adalah $r :
s$ ($r$ banding $s$), dengan $r/s = p/(1-p)$.  Jika $r$ dan $s$ diberikan, maka $p$ 
dapat diperoleh menggunakan persamaan $p = r/(r+s)$.
\end{definition}
\begin{example}(Lanjutan Contoh~\ref{exam 1.7})\label{exam 1.9}
Dalam Contoh~\ref{exam 1.7}, kita menetapkan peluang 1/5 pada kejadian bahwa
kandidat C memenangi pemilihan.  Jadi, rasio peluang yang mendukung kemenangan C adalah $1/5 :
4/5$.  Rasio ini juga dapat ditulis sebagai $1 : 4$, $2
: 8$, dan seterusnya.  Taruhan bahwa C menang bersifat adil jika kita menerima 4 dolar
ketika C menang dan membayar 1 dolar ketika C kalah.
\end{example}

\subsection*{Ruang Sampel Takhingga}\index{ruang sampel!takhingga}

Jika suatu ruang sampel memiliki takhingga banyak titik, cara mendefinisikan
fungsi distribusi
bergantung pada apakah ruang sampel tersebut terhitung.  Suatu
ruang
sampel disebut \emx {takhingga terhitung}\index{ruang sampel!takhingga terhitung} 
jika elemen-elemennya dapat dihitung, yaitu dapat
dipasangkan satu-satu dengan bilangan bulat positif, dan disebut \emx {takhingga
takterhitung} jika tidak.  Secara umum, ruang sampel takhingga memerlukan konsep baru
\choice{}{(lihat Bab~\ref{chp 2})}, 
tetapi ruang takhingga terhitung tidak.  Jika
$$
\Omega = \{\omega_1,\omega_2,\omega_3, \dots\}
$$
adalah ruang sampel takhingga terhitung, fungsi distribusi didefinisikan
tepat seperti dalam Definisi \ref{def 1.2}, kecuali bahwa jumlahnya kini harus
berupa deret takhingga yang \emx {konvergen}. Teorema~\ref{thm 1.1}
tetap berlaku, demikian pula perluasannya, Teorema~\ref{thm 1.1.5} dan \ref{thm 1.2.5}.  
Satu hal yang tidak dapat kita lakukan pada ruang sampel takhingga terhitung, tetapi dapat dilakukan pada 
ruang sampel terhingga, adalah
mendefinisikan fungsi distribusi \emx {seragam} seperti dalam Definisi~\ref{def
1.3}.
Dalam Latihan~\ref{exer 1.2.19}, Anda diminta menjelaskan mengapa hal ini tidak
mungkin.
\begin{example}\label{exam 1.10}
Sebuah koin dilempar sampai sisi kepala muncul untuk pertama kalinya.  Misalkan hasil
percobaan, yaitu $\omega$, adalah nomor pelemparan ketika kepala pertama kali muncul.  Maka
hasil yang mungkin dari percobaan kita adalah
$$
\Omega = \{1,2,3, \dots\}\ .
$$
Perhatikan bahwa meskipun koin dapat saja selalu menunjukkan ekor, kita belum
memasukkan kemungkinan ini.  Alasannya akan segera dijelaskan.  Peluang
kepala muncul pada pelemparan pertama adalah 1/2.  Peluang
ekor muncul pada pelemparan pertama dan kepala pada pelemparan kedua adalah 1/4.  Peluang
memperoleh dua ekor yang diikuti satu kepala adalah 1/8, dan seterusnya. 
Hal ini menyarankan penetapan fungsi distribusi $m(n) = 1/2^n$ untuk $n = 1$,~2, 3, \dots. 
Untuk
melihat bahwa ini merupakan fungsi distribusi, kita harus menunjukkan bahwa
$$
\sum_{\omega} m(\omega) = \frac12 + \frac14 + \frac18 + \cdots = 1\ .
$$
Kebenaran persamaan ini mengikuti rumus jumlah deret geometri\index{deret
geometri},
$$
1 + r + r^2 + r^3 + \cdots = \frac1{1-r}\ ,
$$
atau
\begin{equation}
r + r^2 + r^3 + r^4 + \cdots = \frac r{1-r}\ ,
\label{eq 1.2}
\end{equation} 
untuk $-1 < r < 1$.
\par
Dengan menetapkan $r = 1/2$, kita melihat bahwa peluang koin pada akhirnya
menunjukkan kepala adalah 1.  Hasil yang mungkin berupa ekor pada setiap pelemparan harus
diberi
peluang 0, sehingga kita mengeluarkannya dari ruang sampel
hasil yang mungkin.

Misalkan $E$ adalah kejadian bahwa kepala pertama kali muncul setelah jumlah
pelemparan yang genap.  Maka
$$
E = \{2,4,6,8, \dots\}\ ,
$$
dan
$$
P(E) = \frac14 + \frac1{16} + \frac1{64} +\cdots\ .
$$
Dengan menetapkan $r = 1/4$ dalam Persamaan~\ref{eq 1.2}, kita melihat bahwa              
$$
P(E) = \frac{1/4}{1 - 1/4} = \frac13\ .
$$
Jadi, peluang bahwa kepala pertama kali muncul setelah jumlah
pelemparan yang genap adalah 1/3, sedangkan setelah jumlah pelemparan yang ganjil adalah 2/3.
\end{example}

\subsection*{Catatan Sejarah}

Sebuah pertanyaan menarik dalam sejarah ilmu pengetahuan adalah: Mengapa peluang belum
dikembangkan hingga abad keenam belas?  Kita mengetahui bahwa pada abad keenam belas,
masalah perjudian dan permainan untung-untungan membuat orang mulai memikirkan
peluang.  Namun, perjudian dan permainan untung-untungan hampir sama tuanya dengan
peradaban itu sendiri.  Di Mesir kuno\index{Mesir} (pada masa Dinasti Pertama,
sekitar~3500~{\footnotesize\rm {B.C.}}), dimainkan suatu permainan yang kini disebut ``Hounds and Jackals."
Dalam permainan ini, pergerakan anjing pemburu dan jakal didasarkan pada hasil
pelemparan
dadu bersisi empat yang terbuat dari tulang hewan, yang disebut astragali. 
Dadu bersisi
enam dari berbagai bahan telah ada sejak abad keenam belas
~{\footnotesize\rm {B.C.}}  Perjudian
tersebar luas di Yunani kuno\index{Yunani} dan Roma\index{Roma}.  Bahkan, 
di Kekaisaran Romawi kadang-kadang
dianggap perlu untuk memberlakukan undang-undang yang melarang perjudian.  Lalu, mengapa
peluang belum dihitung hingga abad keenam belas?

Beberapa penjelasan telah diajukan mengenai perkembangan yang terlambat ini.  Salah satunya adalah
bahwa
matematika yang relevan belum dikembangkan dan memang tidak mudah dikembangkan.  Notasi
matematika kuno membuat perhitungan numerik rumit, sedangkan notasi aljabar yang
kita kenal baru dikembangkan pada abad keenam belas. 
Namun, seperti akan kita lihat, banyak gagasan kombinatorial yang diperlukan untuk menghitung
peluang telah dibahas jauh sebelum abad keenam belas.  Karena banyak
kejadian acak pada masa itu berkaitan dengan pengundian yang berhubungan dengan
urusan
keagamaan, diduga mungkin terdapat hambatan agama terhadap
kajian keacakan dan perjudian.  Dugaan lain adalah bahwa diperlukan
pendorong yang lebih kuat, seperti perkembangan perdagangan.  Meskipun demikian, tidak satu pun
penjelasan ini tampak sepenuhnya memuaskan, dan orang masih bertanya-tanya
mengapa diperlukan waktu begitu lama sebelum peluang dikaji secara serius.  Pembahasan yang
menarik mengenai masalah ini dapat ditemukan dalam karya Hacking\index{HACKING, I.}.\footnote{I.~Hacking,
\emx
{The Emergence of Probability} (Cambridge: Cambridge University Press,
1975).}
\par
Orang pertama yang menghitung peluang secara sistematis adalah Gerolamo
Cardano\index{CARDANO, G.|(}
(1501--1576) dalam bukunya \emx {Liber de Ludo Aleae.}  Karya ini diterjemahkan dari
bahasa Latin oleh Gould dan dimuat dalam buku \emx {Cardano: The Gambling
Scholar}
karya Ore\index{ORE, O.}.\footnote{O.~Ore, \emx {Cardano: The Gambling Scholar}
(Princeton:
Princeton University Press, 1953).}  Ore memberikan pembahasan menarik mengenai
kehidupan cendekiawan penuh warna ini, termasuk minatnya pada berbagai
bidang seperti kedokteran, astrologi, dan matematika.  Di sana Anda juga akan
menemukan uraian terperinci mengenai perseteruan Cardano yang terkenal dengan Tartaglia
tentang penyelesaian persamaan kubik.

Dalam bukunya tentang peluang, Cardano hanya membahas kasus khusus yang kita
sebut fungsi distribusi seragam.  Pembatasan pada hasil-hasil berpeluang sama ini
bertahan
untuk waktu yang lama.  Dalam kasus ini, Cardano menyadari bahwa
peluang terjadinya suatu kejadian adalah rasio jumlah hasil yang
mendukung terhadap jumlah seluruh hasil.

Sebagian besar contoh Cardano membahas pelemparan dadu.  Di sini ia menyadari bahwa
hasil
dua pelemparan seharusnya dipandang sebagai 36 pasangan terurut $(i,j)$,
bukan 21 pasangan takterurut.  Ini merupakan pokok yang halus dan masih
menimbulkan masalah jauh kemudian bagi penulis lain tentang peluang.  Sebagai contoh,
pada
abad kedelapan belas, matematikawan Prancis terkenal d'Alembert, penulis beberapa
karya tentang peluang, menyatakan bahwa ketika sebuah koin dilempar dua kali,
jumlah kepala yang muncul adalah 0, 1, atau~2, sehingga kita seharusnya
menetapkan peluang yang sama pada ketiga hasil yang mungkin
tersebut.\footnote{J.~d'Alembert, ``Croix ou Pile," dalam \emx
{L'Encyclop\'edie,}
ed.~Diderot, vol.~4 (Paris, 1754).}  Cardano memilih ruang sampel yang benar
untuk
masalah dadunya dan menghitung peluang yang benar bagi berbagai
kejadian.
\par
Karya matematika Cardano diselingi banyak nasihat bagi calon
penjudi dalam paragraf-paragraf pendek yang, misalnya, berjudul ``Who Should
Play
and When," ``Why Gambling Was Condemned by Aristotle," ``Do Those Who Teach
Also Play Well?," dan sebagainya. Dalam paragraf berjudul ``The Fundamental
Principle of Gambling," Cardano menjadikan kesetaraan kondisi sebagai asas
utama: lawan, penonton, uang, keadaan, kotak dadu, dan dadu itu sendiri harus
diperlakukan setara. Penyimpangan yang menguntungkan lawan merugikan diri
sendiri, sedangkan penyimpangan yang menguntungkan diri sendiri tidak adil.
\footnote{O.~Ore, op.\ cit., p.~189. Ringkasan editorial; ungkapan terjemahan
modern tidak diterjemahkan ulang atau direproduksi.}

Cardano memang membuat kesalahan, dan jika kemudian menyadarinya, ia tidak kembali untuk
memperbaiki kesalahan tersebut.  Sebagai contoh, untuk kejadian yang menguntungkan dalam tiga dari
empat kasus,
Cardano menetapkan rasio peluang yang benar, yaitu $3 : 1$, bahwa kejadian tersebut akan
terjadi.  Namun, ia kemudian menetapkan rasio peluang dengan menguadratkan bilangan-bilangan ini (yaitu $9 :
1$) untuk kejadian yang terjadi dua kali berturut-turut.  Kemudian, dengan mempertimbangkan kasus
ketika rasio peluangnya $1 : 1$, ia menyadari bahwa hal ini tidak mungkin benar dan
sampai pada hasil yang benar: ketika $f$ dari $n$ hasil
menguntungkan, rasio peluang bagi hasil yang menguntungkan dua kali berturut-turut adalah $f^2 : n^2 -
f^2$.  Ore menunjukkan bahwa hal ini ekuivalen dengan pemahaman bahwa jika
peluang suatu kejadian terjadi dalam satu percobaan adalah $p$, maka peluang
kejadian itu
terjadi dua kali adalah $p^2$.  Cardano kemudian membuktikan bahwa untuk tiga
keberhasilan rumusnya adalah $p^3$ dan untuk empat keberhasilan $p^4$, sehingga
jelas bahwa ia memahami bahwa peluangnya adalah $p^n$ untuk $n$ keberhasilan dalam
$n$
pengulangan independen percobaan tersebut.  Hal ini akan mengikuti dari
konsep independensi yang kita perkenalkan dalam Bagian~\ref{sec 4.1}\index{CARDANO, G.|)}. 
\par
Karya Cardano merupakan upaya pertama yang luar biasa untuk menuliskan hukum-hukum
peluang, tetapi bukan karya itulah yang memicu kajian sistematis
tentang bidang ini.  Pemicunya adalah rangkaian surat terkenal antara Pascal dan Fermat. 
Korespondensi ini dimulai oleh Pascal untuk meminta pendapat Fermat mengenai masalah yang
diberikan kepadanya oleh Chevalier de M\'er\'e,\index{de M\'ER\'E, CHEVALIER} seorang penulis terkenal,
tokoh terkemuka di istana Louis XIV, dan penjudi yang bersemangat.
\par
Masalah pertama yang diajukan de M\'er\'e adalah masalah dadu.  Menurut cerita, ia
bertaruh bahwa setidaknya satu angka enam akan muncul dalam empat kali pelemparan sebuah dadu dan
terlalu sering menang, sehingga ia kemudian bertaruh bahwa sepasang angka enam akan muncul dalam 24 kali
pelemparan sepasang dadu.  Peluang munculnya angka enam pada satu dadu adalah 1/6 dan, menurut
hukum perkalian bagi percobaan independen, peluang munculnya dua angka enam ketika
sepasang dadu dilempar adalah $(1/6)(1/6) = 1/36$.  Ore\footnote{O.\ Ore,
``Pascal and the Invention of Probability Theory,'' \emx {American Mathematics
Monthly}, vol. 67 (1960), pp. 409--419.}\index{ORE, O.} menyatakan bahwa suatu aturan perjudian pada
masa itu menyarankan bahwa, karena empat pengulangan menguntungkan bagi terjadinya
kejadian dengan peluang 1/6, bagi kejadian yang enam kali lebih kecil kemungkinannya, $6 \cdot 4 =
24$ pengulangan akan cukup untuk taruhan yang menguntungkan.  Melalui
perhitungan eksak, Pascal menunjukkan bahwa diperlukan 25 pelemparan agar taruhan pada sepasang
angka enam menguntungkan.
\par
Masalah kedua jauh lebih sulit: ini merupakan masalah lama tentang
penentuan pembagian taruhan yang adil dalam sebuah turnamen ketika
rangkaian pertandingan, karena alasan tertentu, dihentikan sebelum selesai.  Masalah ini
kini disebut masalah pembagian taruhan.\index{masalah pembagian taruhan}  Masalah tersebut
telah menjadi masalah baku dalam teks matematika; masalah itu muncul dalam buku Fra Luca
Paccioli, \emx {summa de Arithmetica, Geometria, Proportioni et
Proportionalit\`a}, yang dicetak di Venesia pada 1494.\footnote{ibid., p. 414.
Ringkasan editorial; ungkapan terjemahan modern tidak diterjemahkan ulang atau
direproduksi.} Paccioli membahas permainan bola yang dimenangi pada 60 poin,
dengan 10 poin per babak dan taruhan 10 dukat. Permainan terhenti ketika satu
pihak memiliki 50 poin dan pihak lain 20; masalahnya ialah menentukan pembagian
taruhan yang adil.
\par
Penyelesaian yang masuk akal, seperti membagi taruhan menurut rasio
pertandingan yang dimenangi masing-masing pemain, telah diajukan, tetapi belum ada penyelesaian benar yang
ditemukan pada masa korespondensi Pascal-Fermat.\index{PASCAL, B.|(}\index{FERMAT, P.|(} 
Surat-surat tersebut terutama membahas upaya Pascal dan Fermat menyelesaikan masalah ini.  Blaise
Pascal (1623--1662) adalah anak ajaib yang menerbitkan
risalahnya tentang irisan kerucut pada usia enam belas tahun dan menciptakan mesin
hitung pada usia delapan belas tahun.  Pada masa surat-surat tersebut, pembuktiannya mengenai berat
atmosfer telah menempatkannya di barisan terdepan
fisikawan sezamannya.  Pierre de~Fermat (1601--1665) adalah ahli hukum terpelajar di
Toulouse yang mempelajari matematika pada waktu luangnya.  Ia disebut oleh sebagian orang sebagai pangeran
matematikawan amatir dan salah satu matematikawan murni terbesar sepanjang masa.
\par
Surat-surat tersebut, yang diterjemahkan oleh Maxine Merrington, dimuat dalam karya sejarah peluang
yang menarik oleh Florence David, \emx {Games, Gods and Gambling}.\footnote{F. N.
David, \emx {Games, Gods and Gambling\/} (London: G. Griffin, 1962), p. 230
ff.}\index{DAVID, F.\ N.}  Dalam surat bertanggal Rabu, 29 Juli 1654, Pascal menjelaskan
kepada Fermat sebuah penyelesaian bagi masalah pembagian taruhan. Ringkasan berikut mempertahankan
isi matematis surat tersebut tanpa menerjemahkan atau mereproduksi ungkapan terjemahan modern yang
dikutip oleh David.
\par
Kedua pemain masing-masing mempertaruhkan 32~pistole dalam pertandingan yang ditentukan oleh tiga
kemenangan. Pada keadaan yang dibahas, pemain pertama telah menang dua kali dan pemain kedua sekali.
Jika pemain pertama memenangi pertandingan berikutnya, ia memperoleh seluruh taruhan, yaitu
64~pistole. Jika ia kalah, kedudukan menjadi imbang dua kemenangan sehingga, bila permainan dihentikan,
masing-masing pemain berhak atas 32~pistole. Karena itu pemain pertama sudah pasti memperoleh
32~pistole; sisa 32~pistole bergantung pada satu pertandingan berpeluang seimbang dan harus dibagi dua.
Pembagian yang adil ialah 48~pistole untuk pemain pertama dan 16~pistole untuk pemain kedua.

Argumen Pascal menghasilkan tabel pada Gambar~\ref{fig 1.9} untuk jumlah yang
menjadi hak pemain A pada setiap titik penghentian.
\putfig{3truein}{PSfig1-9}{Tabel Pascal.}{fig 1.9} 
\par
Setiap entri dalam tabel merupakan rata-rata dari bilangan tepat di atas dan di
sebelah kanannya.  Fakta ini, bersama nilai yang diketahui ketika
turnamen selesai, menentukan semua nilai dalam tabel.  Jika pemain A
memenangi pertandingan pertama, ia memerlukan dua kemenangan lagi dan B memerlukan tiga kemenangan untuk
menang; jadi, jika turnamen dibatalkan, A seharusnya menerima 44 pistole.
\par
Surat tempat Fermat menyampaikan penyelesaiannya telah hilang; tetapi
untungnya, Pascal menjelaskan metode Fermat dalam surat bertanggal Senin, 24
Agustus 1654.  Menurut uraian Pascal,\footnote{ibid., p. 239ff.} jika pemain pertama
masih memerlukan dua kemenangan dan pemain kedua tiga kemenangan, paling banyak empat pertandingan
tambahan diperlukan untuk menentukan pemenang. Metode Fermat memperpanjang semua kemungkinan jalannya
permainan hingga panjang yang sama, yaitu empat pertandingan, lalu menghitung banyaknya urutan yang
membuat masing-masing pemain menang dan membagi taruhan menurut perbandingan hitungan tersebut. Ini
merupakan ringkasan editorial; ungkapan terjemahan modern surat itu tidak diterjemahkan atau direproduksi.
\par
Fermat menyadari bahwa berbagai cara permainan dapat berakhir mungkin tidak
berpeluang sama.  Sebagai contoh, jika A memerlukan dua kemenangan lagi dan B memerlukan tiga kemenangan untuk
menang, dua urutan yang mungkin bagi kemenangan A dalam turnamen adalah WLW dan
LWLW.  Kedua urutan ini tidak memiliki peluang terjadi yang sama.  Untuk menghindari
kesulitan ini, Fermat memperpanjang permainan dengan menambahkan pertandingan semu agar semua
cara pertandingan dapat berlangsung memiliki panjang yang sama, yaitu empat.  Ia cukup
cerdik untuk menyadari bahwa perpanjangan ini tidak akan mengubah pemenang dan
bahwa kini ia cukup menghitung jumlah urutan yang menguntungkan masing-masing pemain
karena semuanya telah dibuat berpeluang sama.  Jika kita mendaftarkan semua cara yang mungkin bagi
permainan empat pertandingan yang diperpanjang ini, kita memperoleh 16 hasil permainan yang
mungkin berikut:
\vskip .1in
\begin{table}[h]
\centering
\begin{tabular}{llll}
{\underline{WWWW}} & {\underline{WLWW}} & {\underline{LWWW}} &
{\underline{LLWW}} \\
{\underline{WWWL}} & {\underline{WLWL}} &
{\underline{LWWL}} & {LLWL} \\
{\underline{WWLW}} &
{\underline{WLLW}} & {\underline{LWLW}} & {LLLW} \\
{\underline{WWLL}} & {WLLL} & {LWLL} & {LLLL}\ .
\end{tabular}
\end{table}
\vskip .1in
Pemain A menang dalam kasus yang memuat setidaknya dua kemenangan (11 kasus yang
digarisbawahi), sedangkan B menang dalam kasus yang memuat setidaknya tiga
kekalahan (5 kasus lainnya).  Karena A menang dalam 11 dari 16 kasus yang mungkin, Fermat
berpendapat bahwa peluang A menang adalah 11/16.  Jika taruhannya 64
pistole, A seharusnya menerima 44 pistole, sesuai dengan hasil Pascal. 
Pascal dan Fermat mengembangkan metode yang lebih sistematis untuk menghitung jumlah
hasil yang menguntungkan bagi masalah semacam ini, dan hal tersebut akan menjadi salah satu masalah utama
kita.  Metode penghitungan semacam itu termasuk dalam bidang \emx {kombinatorika},
yang menjadi pokok Bab~\ref{chp 3}.
\par
Kita melihat bahwa kedua matematikawan ini sampai pada dua cara yang sangat berbeda untuk
menyelesaikan masalah pembagian taruhan.  Metode Pascal adalah mengembangkan algoritme dan
menggunakannya untuk menghitung pembagian yang adil.  Metode ini mudah diterapkan pada
komputer dan mudah diperumum.  Sebaliknya, metode Fermat adalah
mengubah masalah menjadi masalah ekuivalen yang dapat ditangani dengan metode pencacahan
atau kombinatorial.  Dalam Bab~\ref{chp 3}, kita akan melihat bahwa Fermat sebenarnya menggunakan 
apa yang kemudian dikenal sebagai segitiga Pascal!  Dalam kajian peluang masa kini, kita 
akan mendapati bahwa pendekatan algoritmik dan pendekatan kombinatorial mendapat tempat yang
setara, sama seperti 300 tahun yang lalu ketika kajian peluang bermula.
\index{PASCAL, B.|)}\index{FERMAT, P.|)}

\exercises
\begin{LJSItem}

\i\label{exer1.2.1} Misalkan $\Omega = \{a,b,c\}$ adalah ruang sampel.  Misalkan $m(a)
= 1/2$, $m(b) =
1/3$, dan $m(c) = 1/6$.  Tentukan peluang bagi kedelapan himpunan bagian dari
$\Omega$.

\i\label{exer 1.2.2} Berikan salah satu ruang sampel yang mungkin, $\Omega$, bagi setiap
percobaan berikut:
\begin{enumerate}
\item Sebuah pemilihan menentukan pemenang antara dua kandidat, A dan B.
\item Sebuah koin bersisi dua dilempar.
\item Seorang mahasiswa ditanya mengenai bulan dan hari dalam sepekan
ketika ulang tahunnya jatuh.
\item Seorang mahasiswa dipilih secara acak dari kelas yang berisi sepuluh mahasiswa.
\item Anda memperoleh suatu nilai dalam mata kuliah ini.
\end{enumerate}

\i\label{exer 1.2.3} Untuk kasus mana dalam Latihan~\ref{exer 1.2.2}  
penggunaan fungsi distribusi seragam merupakan pilihan yang wajar?

\i\label{exer 1.2.4} Jelaskan dengan kata-kata kejadian yang dinyatakan oleh
himpunan bagian berikut dari 
$$\Omega = \{HHH,\ HHT,\ HTH,\ HTT,\ THH,\ THT,\ TTH,\ TTT\}$$
(lihat Contoh~\ref{exam 1.5}).
\begin{enumerate}
\item $E = \{\mbox{HHH,HHT,HTH,HTT}\}$. 
\item $E = \{\mbox{HHH,TTT}\}$.
\item $E = \{\mbox{HHT,HTH,THH}\}$. 
\item $E =
\{\mbox{HHT,HTH,HTT,THH,THT,TTH,TTT}\}$.
\end{enumerate}

\i\label{exer 1.2.5} Berapakah peluang kejadian yang dijelaskan dalam
Latihan~\ref{exer 1.2.4}?

\i\label{exer 1.2.6} Sebuah dadu diberi pemberat sedemikian rupa sehingga peluang
munculnya setiap sisi
sebanding dengan jumlah titik pada sisi tersebut.  (Sebagai contoh, angka enam
tiga kali
lebih mungkin daripada angka dua.)  Berapakah peluang memperoleh bilangan
genap
dalam satu pelemparan?

\i\label{exer 1.2.7} Misalkan $A$ dan $B$ adalah kejadian sedemikian rupa sehingga $P(A \cap B) =
1/4$, $P(\tilde A) =
1/3$, dan $P(B) = 1/2$.  Berapakah $P(A \cup B)$?

\i\label{exer 1.2.8} Seorang mahasiswa harus memilih satu dari mata kuliah seni,
geologi, atau psikologi
sebagai mata kuliah pilihan.  Peluang ia memilih seni sama dengan peluang memilih psikologi, sedangkan peluangnya
memilih geologi dua kali lebih besar.  Secara berurutan, berapakah peluang ia
memilih seni, geologi, dan psikologi?

\i\label{exer 1.2.9} Seorang mahasiswa harus memilih tepat dua dari tiga mata kuliah
pilihan: seni, bahasa Prancis,
dan matematika.  Ia memilih seni dengan peluang 5/8, bahasa Prancis dengan
peluang
5/8, serta seni dan bahasa Prancis sekaligus dengan peluang 1/4.  Berapakah
peluang
ia memilih matematika?  Berapakah peluang ia memilih seni
atau bahasa Prancis?

\i\label{exer 1.2.10} Agar sebuah rancangan undang-undang diajukan kepada Presiden
Amerika Serikat, rancangan itu harus
disahkan oleh Dewan Perwakilan Rakyat dan Senat.  Asumsikan bahwa dari
semua
rancangan undang-undang yang diajukan kepada kedua lembaga ini, 60 persen disahkan oleh Dewan, 80
persen
disahkan oleh Senat, dan 90 persen disahkan oleh setidaknya salah satu dari keduanya.  Hitung
peluang bahwa rancangan undang-undang berikutnya yang diajukan kepada kedua lembaga akan diteruskan kepada
Presiden.

\i\label{exer 1.2.11} Berapakah rasio peluang yang semestinya diberikan seseorang untuk mendukung
kejadian berikut?
\begin{enumerate}
\item Sebuah kartu yang dipilih secara acak dari setumpuk 52 kartu adalah kartu as.
\item Dua kepala akan muncul ketika sebuah koin dilempar dua kali.
\item Boxcars (dua angka enam) akan muncul ketika dua dadu dilempar.
\end{enumerate}

\i\label{exer 1.2.12} Anda menawarkan rasio peluang $3 : 1$ bahwa teman Anda, Smith, akan
terpilih sebagai wali kota
di kota Anda.  Peluang berapakah yang Anda tetapkan bagi kejadian bahwa Smith
menang?

\i\label{exer 1.2.13} Dalam sebuah pacuan kuda, rasio peluang kemenangan Romance
tercantum sebagai $2 :
3$, sedangkan rasio kemenangan Downhill adalah $1 : 2$.  Berapakah rasio peluang yang semestinya diberikan bagi
kejadian bahwa Romance atau Downhill menang?

\i\label{exer 1.2.13.1} Misalkan $X$ adalah peubah acak dengan fungsi distribusi
$m_X(x)$ yang didefinisikan oleh
$$
m_X(-1) = 1/5,\ \  m_X(0) = 1/5,\ \ m_X(1) = 2/5,\ \  m_X(2) = 1/5\ .
$$
\begin{enumerate}
\item Misalkan $Y$ adalah peubah acak yang didefinisikan oleh persamaan $Y = X + 3$.  Tentukan
fungsi distribusi $m_Y(y)$ dari $Y$.
\item Misalkan $Z$ adalah peubah acak yang didefinisikan oleh persamaan $Z = X^2$.  Tentukan 
fungsi distribusi $m_Z(z)$ dari $Z$.
\end{enumerate}

\istar\label{exer 1.2.14} John dan Mary mengikuti mata kuliah matematika.  Mata kuliah ini
hanya memiliki
tiga nilai: A, B, dan C.  Peluang John memperoleh B adalah .3.  Peluang
Mary memperoleh B adalah .4.  Peluang bahwa keduanya tidak memperoleh A,
tetapi setidaknya salah seorang memperoleh B, adalah .1.  Berapakah peluang setidaknya salah seorang
memperoleh B, tetapi keduanya tidak memperoleh C?

\i\label{exer 1.2.15} Dalam sebuah pertempuran sengit, tidak kurang dari 70 persen
prajurit kehilangan satu
mata, tidak kurang dari 75 persen kehilangan satu telinga, tidak kurang dari 80 persen kehilangan satu
tangan, dan tidak kurang dari 85 persen kehilangan satu kaki.  Berapakah persentase minimum yang
mungkin
bagi mereka yang sekaligus kehilangan satu telinga, satu mata, satu tangan, dan
satu
kaki?\footnote{Lihat Knot X, dalam Lewis Carroll, \emx {Mathematical Recreations,}
vol.~2 (Dover, 1958).}

\istar\label{exer 1.2.16}
Asumsikan bahwa peluang sebuah ``keberhasilan" dalam
satu percobaan dengan $n$ hasil adalah $1/n$.  Misalkan $m$ adalah jumlah
percobaan yang diperlukan agar taruhan bahwa setidaknya satu keberhasilan
akan
terjadi menjadi menguntungkan (lihat Latihan~\ref{sec 1.1}.\ref{exer 1.1.5}).

\begin{enumerate}
\item Tunjukkan bahwa peluang tidak ada keberhasilan dalam $m$ uji coba
adalah $(1 - 1/n)^m$.

\item (de Moivre)\index{de MOIVRE, A.} Tunjukkan bahwa jika $m = n \log 2$, maka
$$
\lim_{n \to \infty} \left(1 - \frac1n \right)^m = \frac12\ .
$$
\emx {Petunjuk}:
$$
\lim_{n \to \infty} \left(1 - \frac1n \right)^n = e^{-1}\ .
$$
Jadi, untuk $n$ besar, kita sebaiknya memilih $m$ kira-kira sebesar $n \log 2$.

\item Apakah DeMoivre akan sampai pada jawaban yang benar bagi dua taruhan de~M\'er\'e
\index{de M\'ER\'E, CHEVALIER}
jika ia menggunakan pendekatannya?
\end{enumerate}

\i\label{exer 1.2.17}
\begin{enumerate}
\item Untuk kejadian $A_1$, \dots, $A_n$, buktikan bahwa 
$$P(A_1 \cup \cdots \cup A_n) \leq P(A_1) + \cdots + P(A_n)\ .$$
\item Untuk kejadian $A$ dan $B$, buktikan bahwa
$$
P(A \cap B) \geq P(A) + P(B) - 1.
$$
\end{enumerate}

\i\label{exer 1.2.18}
Jika $A$, $B$, dan~$C$ adalah sebarang tiga kejadian, tunjukkan bahwa
$$\begin{array}{ll}
P(A \cup B \cup C) &= P(A) + P(B) + P(C) \\
                   &\ \  -\, P(A \cap B) - P(B \cap C) - P(C \cap A) \\
                   &\ \  +\, P(A \cap B \cap C)\ .
\end{array}
$$

\i\label{exer 1.2.19} Jelaskan mengapa kita tidak mungkin mendefinisikan
fungsi distribusi seragam (lihat Definisi~\ref{def 1.3}) 
pada ruang sampel takhingga terhitung.  \emx {Petunjuk}: Asumsikan $m(\omega) = a$
untuk semua
$\omega$, dengan $0 \leq a \leq 1$.  Apakah $m(\omega)$ memiliki semua sifat
fungsi
distribusi?

\i\label{exer 1.2.20} Dalam Contoh~\ref{exam 1.10}, tentukan peluang bahwa
koin
menunjukkan kepala untuk pertama kalinya pada pelemparan kesepuluh, kesebelas, atau kedua belas.

\i\label{exer 1.2.21} Sebuah dadu dilempar sampai angka enam muncul untuk
pertama kalinya.  Kita akan
melihat bahwa peluang hal ini terjadi pada pelemparan ke-$n$ adalah
$(5/6)^{n-1}\cdot(1/6)$.  Dengan menggunakan fakta ini, jelaskan ruang sampel takhingga
dan fungsi distribusi yang sesuai bagi percobaan melempar dadu sampai angka enam
muncul
untuk pertama kalinya.  Verifikasikan bahwa untuk fungsi distribusi Anda, $\sum_{\omega} m(\omega)
=
1$.

\i\label{exer 1.2.22} 
Misalkan $\Omega$ adalah ruang sampel
$$
\Omega = \{0,1,2,\dots\}\ ,
$$
dan definisikan fungsi distribusi dengan
$$
m(j) = (1 - r)^j r\ ,
$$
untuk suatu $r$ tetap, $0 < r < 1$, dan untuk $j = 0, 1, 2, \ldots$.
Tunjukkan bahwa ini merupakan fungsi distribusi bagi $\Omega$.

\i\label{exer 1.2.23} Kalender kita\index{kalender} memiliki siklus 400 tahun.  
B.~H.~Brown\index{BROWN, B. H.}
memperhatikan bahwa
jumlah kemunculan tanggal tiga belas pada setiap hari dalam
sepekan selama
4800 bulan dalam satu siklus adalah sebagai berikut:

\par Minggu 687
\par Senin 685
\par Selasa 685
\par Rabu 687
\par Kamis 684
\par Jumat 688
\par Sabtu 684

\noindent Dari sini ia menyimpulkan bahwa tanggal tiga belas lebih mungkin jatuh pada
hari Jumat daripada hari lainnya.  Jelaskan maksud pernyataan tersebut.

\i\label{exer 1.2.24} Tversky dan Kahneman\footnote{K.~McKean, ``Decisions,
Decisions," pp.~22--31.} meminta sekelompok subjek mengerjakan sebuah tugas yang kini dikenal
sebagai \emx{masalah Linda}.\index{masalah Linda} Profil tokoh fiktif berusia 31 tahun dalam
tugas itu dirancang agar aktivitas dalam gerakan feminis terasa sangat representatif. Alih-alih
mereproduksi profil tersebut, nyatakan
$F=$ kejadian bahwa tokoh itu aktif dalam gerakan feminis dan $B=$ kejadian bahwa ia adalah teller bank.
Para subjek diminta mengurutkan kemungkinan:
\hfill\break
\indent(1)~$F$.\hfill\break
\indent(2)~$B$.\hfill\break
\indent(3)~$B\cap F$.\hfill\break

\noindent Tversky\index{TVERSKY, A.} dan Kahneman\index{KAHNEMAN, D.} menemukan bahwa 
antara 85 dan 90 persen subjek menilai
alternatif (1) sebagai yang paling mungkin, tetapi menilai alternatif (3) lebih mungkin daripada alternatif
(2).  Benarkah demikian?  Mereka menyebut fenomena ini \emx {kekeliruan konjungsi},\index{kekeliruan
konjungsi}\index{kekeliruan} dan
mencatat bahwa fenomena ini tampaknya tidak dipengaruhi oleh pelatihan sebelumnya dalam peluang atau
statistika.  Apakah fenomena ini merupakan kekeliruan?  Jika ya, mengapa? Dapatkah Anda memberikan
penjelasan yang mungkin bagi pilihan para subjek?

\i\label{exer 1.2.25} Dua kartu diambil secara berurutan dari setumpuk 52
kartu.  Tentukan
peluang bahwa peringkat kartu kedua lebih tinggi daripada kartu pertama.  \emx
{Petunjuk}: Tunjukkan bahwa $1 = P(\mbox{lebih tinggi}) + P(\mbox{lebih rendah}) + P(\mbox{sama})$
dan gunakan fakta bahwa $P(\mbox{lebih tinggi}) = P(\mbox{lebih rendah})$.

\i\label{exer 1.2.26} \emx {Tabel hayat}\index{tabel hayat} adalah tabel yang, untuk suatu
jumlah kelahiran tertentu, mencantumkan perkiraan jumlah orang yang akan hidup sampai
usia tertentu.  Dalam Lampiran~C, kami memberikan tabel hayat berdasarkan 100{,}000
kelahiran untuk usia 0 sampai 85, baik bagi perempuan maupun laki-laki.  Tunjukkan bagaimana 
tabel ini dapat digunakan untuk memperkirakan peluang $m(x)$ bahwa seseorang yang lahir pada 1981 
akan hidup sampai usia $x$.  Tulis program untuk menggambarkan $m(x)$ bagi laki-laki dan perempuan, 
lalu komentari perbedaan yang Anda lihat antara kedua kasus.

\istar\label{exer 1.2.27} Berikut adalah upaya untuk mengatasi kenyataan bahwa kita
tidak dapat 
memilih suatu ``bilangan bulat acak."\index{bilangan bulat acak}
\begin{enumerate}
\item Secara intuitif, berapakah peluang bahwa bilangan bulat positif yang
``dipilih secara acak" merupakan kelipatan 3?

\item Misalkan $P_3(N)$ adalah peluang bahwa suatu bilangan bulat yang dipilih secara acak
antara 1 dan $N$ merupakan kelipatan 3 (karena ruang sampelnya terhingga, ini
merupakan peluang yang sah).  Tunjukkan bahwa limit
$$
P_3 = \lim_{N \to \infty} P_3(N)
$$
ada dan sama dengan 1/3.  Hal ini memformalkan intuisi pada
(a), dan memberi kita cara untuk menetapkan ``peluang" pada
kejadian tertentu yang merupakan himpunan bagian takhingga dari bilangan bulat positif.

\item Jika $A$ adalah sebarang himpunan bilangan bulat positif, misalkan $A(N)$ menyatakan jumlah
elemen $A$ yang kurang dari atau sama dengan $N$.  Kemudian definisikan
``peluang" $A$ sebagai 
$$
P(A) = \lim_{N \to \infty} A(N)/N\ ,
$$
asalkan limit ini ada.  Tunjukkan bahwa definisi ini akan menetapkan
peluang 0 pada setiap himpunan terhingga dan peluang 1 pada himpunan semua bilangan bulat positif. 
Jadi, peluang himpunan semua bilangan bulat bukanlah jumlah dari
peluang setiap bilangan bulat dalam himpunan tersebut.  Ini berarti bahwa
definisi peluang yang diberikan di sini belum sepenuhnya
memuaskan.

\item Misalkan $A$ adalah himpunan semua bilangan bulat positif dengan jumlah digit ganjil.
Tunjukkan bahwa $P(A)$ tidak ada.  Hal ini menunjukkan bahwa menurut definisi
peluang di atas, tidak semua himpunan memiliki peluang.

\end{enumerate}


\i\label{exer 1.2.28}
(dari Sholander\footnote{M. Sholander, Problem \#1034, \emx {Mathematics Magazine,} vol.~52,
no.\ 3 (May 1979), p.~183.})\index{SHOLANDER, M.}\index{simpang susun semanggi} Dalam sebuah
simpang susun semanggi baku, terdapat empat jalur landai untuk berbelok ke kanan, dan di sebelah dalam
keempat jalur tersebut terdapat empat jalur landai lain untuk berbelok ke kiri.  Mobil Anda mendekati
simpang susun dari arah selatan.  Sebuah mekanisme dipasang sehingga pada setiap titik yang
menawarkan pilihan arah, mobil berbelok ke kanan dengan peluang tetap~$r$. 
\begin{enumerate}
\item
Jika $r = 1/2$, berapakah peluang Anda keluar dari simpang susun menuju barat?
\item
Tentukan nilai $r$ yang memaksimumkan peluang Anda keluar dari simpang susun menuju barat.
\end{enumerate}

\i\label{exer 1.2.29}
(dari Benkoski\footnote{S. Benkoski, Comment on Problem \#1034, \emx {Mathematics Magazine,}
vol.~52, no.\ 3 (May 1979), pp.~183-184.})\index{BENKOSKI, S.} Pertimbangkan simpang susun
semanggi ``murni" yang tidak memiliki jalur landai untuk berbelok ke kanan, melainkan hanya dua
jalan raya lurus yang berpotongan dengan jalur semanggi untuk berbelok ke kiri.  (Jadi, untuk berbelok ke kanan di
simpang susun semacam ini, seseorang harus berbelok ke kiri tiga kali.)  Seperti dalam masalah sebelumnya, mobil Anda
mendekati simpang susun dari arah selatan.  Berapakah nilai $r$ yang memaksimumkan
peluang Anda keluar dari simpang susun menuju timur?

\i\label{exer 1.2.30}
(berdasarkan pertanyaan yang dibahas vos~Savant\footnote{M.~vos~Savant,
\emx{Parade Magazine}, 3 March 1996, p. 14.}\index{vos SAVANT, M.}) Dalam ujian
susulan, dua mahasiswa terlebih dahulu menjawab pertanyaan 5 poin, kemudian ditempatkan
terpisah untuk menjawab pertanyaan 95 poin yang meminta mereka menyebut salah satu dari empat
posisi ban mobil. Jika keduanya memilih posisi secara acak dan independen, berapakah peluang
bahwa jawaban mereka sama? Ringkasan mandiri ini mempertahankan persoalan peluangnya tanpa
menerjemahkan atau mereproduksi cerita surat pembaca.


\begin{enumerate}
\item
Apakah jawabannya 1/16?
\item
Pertanyaan berikut diajukan kepada satu kelas mahasiswa.  ``Hari ini saya berkendara ke
kampus dan salah satu ban saya kempis.  Menurut Anda, ban yang mana?"
Jawabannya adalah sebagai berikut: kanan depan, 58\%, kiri depan, 11\%, kanan belakang, 
18\%, kiri belakang, 13\%.  Misalkan distribusi ini berlaku pada populasi umum,
dan asumsikan bahwa kedua peserta ujian dipilih secara acak dari populasi umum.  Berapakah
peluang mereka memberikan jawaban yang sama untuk pertanyaan kedua?
\end{enumerate}


\end{LJSItem}
